\documentclass[11pt,a4paper]{article}
\usepackage[utf8]{inputenc}
\usepackage[T1]{fontenc}
\usepackage[ngerman]{babel}
\usepackage{amsmath,amssymb}
\usepackage[a4paper,margin=25mm]{geometry}
% longtable und tabularx in einem: Tabellen laufen ueber Seiten, und
% Textspalten, die umbrechen duerfen, nehmen die Zeilenbreite ein.
\usepackage{xltabular}
\usepackage{xcolor}
\usepackage{textcomp}
\usepackage{fancyhdr}

\newcommand{\normref}[1]{{\small\textcolor{gray}{#1}}}
\newcommand{\hinweis}[2]{\par\noindent\textcolor{gray}{\textit{#1:}} \textit{#2}\par}

\pagestyle{fancy}
\fancyhf{}
\fancyhead[L]{\small\berichttitel}
\fancyfoot[C]{\thepage}
\renewcommand{\headrulewidth}{0.4pt}
\setlength{\parindent}{0pt}
\setlength{\parskip}{0.6ex}

\newcommand{\berichttitel}{Altformat}
\title{Altformat}
\author{}
\date{\today}
\begin{document}
\maketitle
40 Berechnungen ausgeführt, 87 Werte bestimmt.

\section{Herleitung}
\subsection{Querschnittsanalyse: Unterzug}
\noindent\textbf{Fläche und Schwerpunkt eines Polygons (Satz von Gauss)}
\begin{equation*}
A = \frac{1}{2} \sum_i c_i \qquad y_S = \frac{1}{6A} \sum_i (y_i + y_{i+1})\, c_i \qquad z_S = \frac{1}{6A} \sum_i (z_i + z_{i+1})\, c_i \qquad c_i = y_i\, z_{i+1} - y_{i+1}\, z_i
\end{equation*}
\noindent\textbf{Trägheitsmomente eines Polygons, bezogen auf den Schwerpunkt}
\begin{equation*}
I_y = \frac{1}{12} \sum_i (z_i^2 + z_i z_{i+1} + z_{i+1}^2)\, c_i - A\, z_S^2\qquad I_z = \frac{1}{12} \sum_i (y_i^2 + y_i y_{i+1} + y_{i+1}^2)\, c_i - A\, y_S^2
\end{equation*}
\noindent\textbf{Polygone des Querschnitts}
\begin{longtable}{lllrrrr}
\hline
Polygon & Material & liegt in & Ecken & $\pm A\ [\mathrm{mm}^2]$ & $y_S\ [\mathrm{mm}]$ & $z_S\ [\mathrm{mm}]$ \\
\hline
\endhead
Polygon 1 & C30/37 & – & $4$ & $180000$ & $150.0$ & $300.0$ \\
\hline
\end{longtable}
\noindent\textbf{Fläche des Bruttoquerschnitts}
\begin{equation*}
A = 180000\,\mathrm{mm}^{2}
\end{equation*}
\noindent\textbf{Schwerpunkt, y}
\begin{equation*}
y_S = 150\,\mathrm{mm}
\end{equation*}
\noindent\textbf{Schwerpunkt, z}
\begin{equation*}
z_S = 300\,\mathrm{mm}
\end{equation*}
\noindent\textbf{Trägheitsmoment um y}
\begin{equation*}
I_y = 5400\,\cdot 10^{6}\,\mathrm{mm}^{4}
\end{equation*}
\noindent\textbf{Trägheitsmoment um z}
\begin{equation*}
I_z = 1350\,\cdot 10^{6}\,\mathrm{mm}^{4}
\end{equation*}
\noindent\textbf{Deviationsmoment}
\begin{equation*}
I_{yz} = 0\,\cdot 10^{6}\,\mathrm{mm}^{4}
\end{equation*}
\noindent\textbf{Teilung einer Stablinie: gleiche Felder, auf die Länge gerundet}
\begin{equation*}
s = \frac{L}{m} \qquad m = \mathrm{runde}\left(\frac{L}{s_{gewählt}}\right)
\end{equation*}
\noindent\textbf{Bewehrung}
\begin{longtable}{lllrr}
\hline
Element & $\text{Lage}\ [\mathrm{mm}]$ & Bewehrung & Stahl & $A_s\ [\mathrm{mm}^2]$ \\
\hline
\endhead
Linie 1 & $(50;\ 50) \rightarrow (250;\ 50)$ & $3 \varnothing 20\ @\ 100.0$ & B500B & $942$ \\
Linie 2 & $(50;\ 550) \rightarrow (250;\ 550)$ & $2 \varnothing 12\ @\ 200.0$ & B500B & $226$ \\
\hline
\end{longtable}
\noindent\textbf{Stahlfläche der Bewehrung}
\begin{equation*}
A_{s,tot} = 1169\,\mathrm{mm}^{2}
\end{equation*}
\noindent\textbf{Bügelquerschnitt je Länge einer Schubwand (n Schnitte)}
\begin{equation*}
\frac{A_{sw}}{s} = \frac{n \cdot \pi\, \varnothing^2 / 4}{s}
\end{equation*}
\noindent\textbf{Schubwände}
\begin{longtable}{llrrlr}
\hline
Wand & $\text{Achse}\ [\mathrm{mm}]$ & $l\ [\mathrm{mm}]$ & $b_w\ [\mathrm{mm}]$ & Bügel & $A_{sw}/s\ [\mathrm{mm}^2/\mathrm{m}]$ \\
\hline
\endhead
Wand 1 & $(50;\ 50) \rightarrow (250;\ 50)$ & $200$ & $100$ & $1 \times \varnothing 10\ @\ 150$ & $524$ \\
Wand 2 & $(250;\ 50) \rightarrow (250;\ 550)$ & $500$ & $100$ & $1 \times \varnothing 10\ @\ 150$ & $524$ \\
Wand 3 & $(250;\ 550) \rightarrow (50;\ 550)$ & $200$ & $100$ & $1 \times \varnothing 10\ @\ 150$ & $524$ \\
Wand 4 & $(50;\ 550) \rightarrow (50;\ 50)$ & $500$ & $100$ & $1 \times \varnothing 10\ @\ 150$ & $524$ \\
\hline
\end{longtable}
\noindent\textbf{Dehnungsebene und innere Kräfte}
\begin{equation*}
\varepsilon(y, z) = \varepsilon_0 + \kappa \cdot v \qquad v = n_y (y - y_S) + n_z (z - z_S) \qquad N = \int_A \sigma\,\mathrm{d}A \qquad M_y = -\int_A \sigma\,(z - z_S)\,\mathrm{d}A \qquad M_z = -\int_A \sigma\,(y - y_S)\,\mathrm{d}A
\end{equation*}
\noindent\textbf{Erfüllungsgrad bei fester Normalkraft und fester Richtung}
\begin{equation*}
\alpha_{eff} = \frac{M_{Rd}}{M_{Ed}} \qquad M = \sqrt{M_y^2 + M_z^2}
\end{equation*}
\noindent\textbf{C30/37: Parabel-Rechteck-Beziehung, Bemessungswerte} \hfill \normref{SIA 262:2025, 4.2.1.6}
\begin{equation*}
\sigma_c(\varepsilon_c) = \begin{cases} -f_{cd} \cdot \dfrac{k_\sigma\,\eta - \eta^2}{1 + (k_\sigma - 2)\,\eta} & 0 \le |\varepsilon_c| \le \varepsilon_{c1d},\ \eta = \dfrac{|\varepsilon_c|}{\varepsilon_{c1d}} \\[2ex] -f_{cd} & \varepsilon_{c1d} < |\varepsilon_c| \le \varepsilon_{c2d} \\[1ex] 0 & \varepsilon_c > 0 \quad (\text{Zug, gerissen}) \end{cases}
\end{equation*}
\noindent\textbf{Werte C30/37}
\begin{longtable}{lr}
\hline
Kennwert & Wert \\
\hline
\endhead
$f_{cd}$ & $20\,\mathrm{N}/\mathrm{mm}^{2}$ \\
$\varepsilon_{c1d}$ & $2\,\text{‰}$ \\
$\varepsilon_{c2d}$ & $3.5\,\text{‰}$ \\
\hline
\end{longtable}
\noindent\textbf{B500B: bilineare Beziehung, Bemessungswerte} \hfill \normref{SIA 262:2025, 4.2.2.4}
\begin{equation*}
\sigma_s(\varepsilon_s) = \begin{cases} \min(E_s\,\varepsilon_s;\ f_{yd}) & \varepsilon_s \ge 0 \\[1ex] \max(E_s\,\varepsilon_s;\ -f_{yd}^{-}) & \varepsilon_s < 0 \end{cases} \qquad |\varepsilon_s| \le \varepsilon_{ud}
\end{equation*}
\noindent\textbf{Werte B500B}
\begin{longtable}{lr}
\hline
Kennwert & Wert \\
\hline
\endhead
$f_{yd}$ & $435\,\mathrm{N}/\mathrm{mm}^{2}$ \\
$f_{yd}^{-}$ & $435\,\mathrm{N}/\mathrm{mm}^{2}$ \\
$E_s$ & $200000\,\mathrm{N}/\mathrm{mm}^{2}$ \\
$\varepsilon_{ud}$ & $4.5\,\%$ \\
\hline
\end{longtable}
\noindent\textbf{Reine Normalkraft}
\begin{longtable}{lr}
\hline
Grenze & Wert \\
\hline
\endhead
reiner Zug, gleichmässig & $N_{Rd}^{+} = 508.1\,\mathrm{kN}$ \\
reiner Druck, gleichmässig & $N_{Rd}^{-} = -4044.1\,\mathrm{kN}$ \\
\hline
\end{longtable}
\subsubsection{Feld}
\noindent\textbf{Einwirkung}
\begin{longtable}{lrl}
\hline
Einwirkung & Wert &  \\
\hline
\endhead
N\_Ed & $0.0$ & kN \\
M\_y,Ed & $150.0$ & kNm \\
M\_z,Ed & $0.0$ & kNm \\
\hline
\end{longtable}
Bruchzustand: Zugrichtung n = (0.000; -1.000), -90.0° gegen die y-Achse. Die Nulllinie steht senkrecht zur Zugrichtung, 229.4 mm auf der Druckseite des Schwerpunkts.

\noindent\textbf{Dehnungsebene}
\begin{longtable}{rr}
\hline
$\varepsilon_0\ [\text{‰}]$ & $\kappa\ [\mathrm{1/m}]$ \\
\hline
\endhead
$11.371$ & $0.04957$ \\
\hline
\end{longtable}
\noindent\textbf{Kräfte im Bruchzustand -- Beton als Resultierende, Bewehrung je Element}
\begin{longtable}{lrrrrr}
\hline
Teil & $F\ [\mathrm{kN}]$ & $y\ [\mathrm{mm}]$ & $z\ [\mathrm{mm}]$ & $\varepsilon\ [\text{‰}]$ & $\sigma\ [\mathrm{N/mm^2}]$ \\
\hline
\endhead
C30/37 & $-367.6$ & $150.0$ & $568.8$ & – & – \\
Linie 1 & $409.8$ & $150.0$ & $50.0$ & $23.76$ & $435$ \\
Linie 2 & $-42.2$ & $150.0$ & $550.0$ & $-1.02$ & $-204$ \\
\hline
\end{longtable}
\noindent\textbf{Probe: Normalkraft}
\begin{equation*}
N_{Rd} = \sum F_i = 0\,\mathrm{kN}
\end{equation*}
\noindent\textbf{Probe: Moment um y}
\begin{equation*}
M_{y,Rd} = -\sum F_i \cdot (z_i - z_S) = 211.8\,\mathrm{kNm}
\end{equation*}
\noindent\textbf{Probe: Moment um z}
\begin{equation*}
M_{z,Rd} = -\sum F_i \cdot (y_i - y_S) = 0\,\mathrm{kNm}
\end{equation*}
\noindent\textbf{Widerstand in der Richtung der Einwirkung}
\begin{equation*}
M_{Rd} = \sqrt{M_{y,Rd}^{2} + M_{z,Rd}^{2}} = \sqrt{\left(211.8\,\mathrm{kNm}\right)^{2} + \left(0\,\mathrm{kNm}\right)^{2}} = 211.8\,\mathrm{kNm}
\end{equation*}
\noindent\textbf{Erfüllungsgrad}
\begin{equation*}
\alpha_{eff,\text{Feld}} = \frac{M_{Rd}}{M_{Ed}} = \frac{211.8\,\mathrm{kNm}}{150\,\mathrm{kNm}} = 1.41 \quad \Rightarrow \quad \text{erfüllt}
\end{equation*}
\noindent\textbf{Querkraft über die Federn der Wände}
\begin{equation*}
k_i = b_{w,i} \cdot l_i \qquad \sum_i k_i\, e_i e_i^T\, u = V \qquad V_i = k_i\, e_i \cdot u \qquad \Delta T = \sum_i V_i\, r_i
\end{equation*}
\noindent\textbf{Torsion in den Zellen (Bredt, mehrzellig)}
\begin{equation*}
T - \Delta T = \sum_k 2 A_k\, q_k \qquad \oint_k \frac{q}{b_w}\, \mathrm{d}s = 2 A_k\, \theta
\end{equation*}
\noindent\textbf{Widerstand je Länge einer Wand (Fachwerk)} \hfill \normref{SIA 262:2025, 4.3.3.4}
\begin{equation*}
v_{Rd,s} = \frac{A_{sw}}{s}\, f_{sd} \cot\alpha \qquad v_{Rd,c} = b_w\, k_c\, f_{cd} \sin\alpha \cos\alpha \qquad v_{Rd} = \max_\alpha \min(v_{Rd,s};\ v_{Rd,c})
\end{equation*}
\noindent\textbf{Druckfeld}
\begin{longtable}{lr}
\hline
Grenze & Wert \\
\hline
\endhead
$\alpha_{min}$ & $30{}^{\circ}$ \\
$\alpha_{max}$ & $45{}^{\circ}$ \\
$k_c$ & $0.55$ \\
\hline
\end{longtable}
\noindent\textbf{Schubwände und ihre Steifigkeit}
\begin{longtable}{lrrrrl}
\hline
Wand & $l\ [\mathrm{mm}]$ & $b_w\ [\mathrm{mm}]$ & $k = b_w l\ [\mathrm{mm}^2]$ & $A_{sw}/s\ [\mathrm{mm^2/m}]$ & Zelle \\
\hline
\endhead
Wand 1 & $200$ & $100$ & $20000$ & $524$ & 1 \\
Wand 2 & $500$ & $100$ & $50000$ & $524$ & 1 \\
Wand 3 & $200$ & $100$ & $20000$ & $524$ & 1 \\
Wand 4 & $500$ & $100$ & $50000$ & $524$ & 1 \\
\hline
\end{longtable}
\subsubsection{Feld}
\noindent\textbf{Einwirkung}
\begin{longtable}{lrl}
\hline
Einwirkung & Wert &  \\
\hline
\endhead
V\_y,Ed & $0.0$ & kN \\
V\_z,Ed & $150.0$ & kN \\
T\_Ed & $15.0$ & kNm \\
\hline
\end{longtable}
\noindent\textbf{Versatz: Moment der Querkraftanteile um den Schwerpunkt}
\begin{equation*}
\Delta T = \sum_i V_i \cdot r_i = 0\,\mathrm{kNm}
\end{equation*}
\noindent\textbf{Was die Zellen tragen}
\begin{equation*}
T_{Zellen} = T_{Ed} - \Delta T = 15\,\mathrm{kNm} - 0\,\mathrm{kNm} = 15\,\mathrm{kNm}
\end{equation*}
\noindent\textbf{Umlauf-Schubfluss der Zellen (Bredt)}
\begin{longtable}{lrr}
\hline
Zelle & $A_k\ [\mathrm{mm}^2]$ & $q_k\ [\mathrm{kN/m}]$ \\
\hline
\endhead
Zelle 1 & $100000$ & $75.0$ \\
\hline
\end{longtable}
\noindent\textbf{Je Wand: Kraft aus der Querkraft, stärkster Schubfluss, Widerstand}
\begin{longtable}{lrrrrrrr}
\hline
Wand & $V_i\ [\mathrm{kN}]$ & $|q_{Ed}|\ [\mathrm{kN/m}]$ & $\alpha\ [{}^{\circ}]$ & $v_{Rd,s}\ [\mathrm{kN/m}]$ & $v_{Rd,c}\ [\mathrm{kN/m}]$ & $v_{Rd}\ [\mathrm{kN/m}]$ & $\alpha_{eff}$ \\
\hline
\endhead
Wand 1 & $0.0$ & $75.0$ & $30$ & $394.3$ & $476.3$ & $394.3$ & $5.26$ \\
Wand 2 & $75.0$ & $225.0$ & $30$ & $394.3$ & $476.3$ & $394.3$ & $1.75$ \\
Wand 3 & $0.0$ & $75.0$ & $30$ & $394.3$ & $476.3$ & $394.3$ & $5.26$ \\
Wand 4 & $-75.0$ & $75.0$ & $30$ & $394.3$ & $476.3$ & $394.3$ & $5.26$ \\
\hline
\end{longtable}
Massgebend: Wand 2.

\noindent\textbf{Erfüllungsgrad}
\begin{equation*}
\alpha_{eff,V,\text{Feld}} = \frac{v_{Rd}}{q_{Ed}} = \frac{394.3\,\mathrm{kN}/\mathrm{m}}{225\,\mathrm{kN}/\mathrm{m}} = 1.75 \quad \Rightarrow \quad \text{erfüllt}
\end{equation*}
\subsection{Beton: C30/37}
\noindent\textbf{Charakteristische Zylinderdruckfestigkeit} \hfill \normref{SIA 262:2025, 3.1.2.2.7}
\begin{equation*}
f_{ck} = 30\,\mathrm{N}/\mathrm{mm}^{2}
\end{equation*}
\noindent\textbf{Beiwert zur Berücksichtigung der Festigkeitsminderung} \hfill \normref{SIA 262:2025, 2.4.2.3}
\begin{equation*}
\eta_{fc} = \min\left[\left(\frac{40\,\mathrm{N}/\mathrm{mm}^{2}}{f_{ck}}\right)^{1/3};\ 1.0\right] = \min\left[\left(\frac{40\,\mathrm{N}/\mathrm{mm}^{2}}{30\,\mathrm{N}/\mathrm{mm}^{2}}\right)^{1/3};\ 1.0\right] = 1
\end{equation*}
\noindent\textbf{Teilsicherheitsbeiwert für Beton} \hfill \normref{SIA 262:2025, 2.4.2.6}
\begin{equation*}
\gamma_c = 1.5
\end{equation*}
\noindent\textbf{Bemessungswert der Betondruckfestigkeit} \hfill \normref{SIA 262:2025, 2.4.2.3}
\begin{equation*}
f_{cd} = \frac{\eta_{fc} \cdot f_{ck}}{\gamma_c} = \frac{1 \cdot 30\,\mathrm{N}/\mathrm{mm}^{2}}{1.5} = 20\,\mathrm{N}/\mathrm{mm}^{2}
\end{equation*}
\noindent\textbf{Dehnung am Ende des ansteigenden Astes} \hfill \normref{SIA 262:2025, 4.2.1.4}
\begin{equation*}
\varepsilon_{c1d} = 2\,\text{‰}
\end{equation*}
\noindent\textbf{Bruchdehnung des Betons} \hfill \normref{SIA 262:2025, 4.2.1.4}
\begin{equation*}
\varepsilon_{c2d} = 3.5\,\text{‰}
\end{equation*}
\noindent\textbf{Beiwert für die Gesteinskörnung} \hfill \normref{SIA 262:2025, 3.1.2.3.3}
\begin{equation*}
k_e = 10000
\end{equation*}
\noindent\textbf{Mittelwert der Zylinderdruckfestigkeit} \hfill \normref{SIA 262:2025, 3.1.2.2.2}
\begin{equation*}
f_{cm} = f_{ck} + 8\,\mathrm{N}/\mathrm{mm}^{2} = 30\,\mathrm{N}/\mathrm{mm}^{2} + 8\,\mathrm{N}/\mathrm{mm}^{2} = 38\,\mathrm{N}/\mathrm{mm}^{2}
\end{equation*}
\noindent\textbf{Mittelwert des Elastizitätsmoduls} \hfill \normref{SIA 262:2025, 3.1.2.3.3}
\begin{equation*}
E_{cm} = k_e \cdot \sqrt[3]{f_{cm}} = 10000 \cdot \sqrt[3]{38} = 33620\,\mathrm{N}/\mathrm{mm}^{2} \quad \left(f_{cm}\ \text{in}\ \mathrm{N}/\mathrm{mm}^{2}\right)
\end{equation*}
\noindent\textbf{Teilsicherheitsbeiwert für den Elastizitätsmodul} \hfill \normref{SIA 262:2025, 4.2.1.15}
\begin{equation*}
\gamma_{cE} = 1
\end{equation*}
\noindent\textbf{Bemessungswert des Elastizitätsmoduls} \hfill \normref{SIA 262:2025, 4.2.1.15}
\begin{equation*}
E_{cd} = \frac{E_{cm}}{\gamma_{cE}} = \frac{33620\,\mathrm{N}/\mathrm{mm}^{2}}{1} = 33620\,\mathrm{N}/\mathrm{mm}^{2}
\end{equation*}
\noindent\textbf{Krümmungsbeiwert der Spannungs-Dehnungs-Beziehung} \hfill \normref{SIA 262:2025, 4.2.1.6}
\begin{equation*}
k_{\sigma} = \frac{E_{cd}}{400 \cdot f_{cd}} = \frac{33620\,\mathrm{N}/\mathrm{mm}^{2}}{400 \cdot 20\,\mathrm{N}/\mathrm{mm}^{2}} = 4.202
\end{equation*}
\noindent\textbf{Mittelwert der Zugfestigkeit} \hfill \normref{SIA 262:2025, 3.1.2.2.7}
\begin{equation*}
f_{ctm} = 2.9\,\mathrm{N}/\mathrm{mm}^{2}
\end{equation*}
\subsection{Betonstahl: B500B}
\noindent\textbf{Charakteristische Fliessgrenze} \hfill \normref{SIA 262:2025, 3.2.2.3}
\begin{equation*}
f_{yk} = 500\,\mathrm{N}/\mathrm{mm}^{2}
\end{equation*}
\noindent\textbf{Teilsicherheitsbeiwert für Betonstahl} \hfill \normref{SIA 262:2025, 2.4.2.6}
\begin{equation*}
\gamma_s = 1.15
\end{equation*}
\noindent\textbf{Bemessungswert der Fliessgrenze} \hfill \normref{SIA 262:2025, 2.4.2.5}
\begin{equation*}
f_{yd} = \frac{f_{yk}}{\gamma_s} = \frac{500\,\mathrm{N}/\mathrm{mm}^{2}}{1.15} = 435\,\mathrm{N}/\mathrm{mm}^{2}
\end{equation*}
\noindent\textbf{Charakteristische Fliessgrenze auf Druck} \hfill \normref{SIA 262:2025, 3.2.2.3}
\begin{equation*}
f_{yk}^{-} = 500\,\mathrm{N}/\mathrm{mm}^{2}
\end{equation*}
\noindent\textbf{Bemessungswert der Fliessgrenze auf Druck} \hfill \normref{SIA 262:2025, 2.4.2.5}
\begin{equation*}
f_{yd}^{-} = \frac{f_{yk}^{-}}{\gamma_s} = \frac{500\,\mathrm{N}/\mathrm{mm}^{2}}{1.15} = 435\,\mathrm{N}/\mathrm{mm}^{2}
\end{equation*}
\noindent\textbf{Elastizitätsmodul des Betonstahls} \hfill \normref{SIA 262:2025, 3.2.2.4}
\begin{equation*}
E_s = 200000\,\mathrm{N}/\mathrm{mm}^{2}
\end{equation*}
\noindent\textbf{Bemessungswert der Dehnung bei Höchstlast} \hfill \normref{SIA 262:2025, 4.2.2.1}
\begin{equation*}
\varepsilon_{ud} = 4.5\,\%
\end{equation*}
\subsection{Querschnittsanalyse: Kastenträger}
\noindent\textbf{Fläche und Schwerpunkt eines Polygons (Satz von Gauss)}
\begin{equation*}
A = \frac{1}{2} \sum_i c_i \qquad y_S = \frac{1}{6A} \sum_i (y_i + y_{i+1})\, c_i \qquad z_S = \frac{1}{6A} \sum_i (z_i + z_{i+1})\, c_i \qquad c_i = y_i\, z_{i+1} - y_{i+1}\, z_i
\end{equation*}
\noindent\textbf{Trägheitsmomente eines Polygons, bezogen auf den Schwerpunkt}
\begin{equation*}
I_y = \frac{1}{12} \sum_i (z_i^2 + z_i z_{i+1} + z_{i+1}^2)\, c_i - A\, z_S^2\qquad I_z = \frac{1}{12} \sum_i (y_i^2 + y_i y_{i+1} + y_{i+1}^2)\, c_i - A\, y_S^2
\end{equation*}
\noindent\textbf{Polygone des Querschnitts}
\begin{longtable}{lllrrrr}
\hline
Polygon & Material & liegt in & Ecken & $\pm A\ [\mathrm{mm}^2]$ & $y_S\ [\mathrm{mm}]$ & $z_S\ [\mathrm{mm}]$ \\
\hline
\endhead
Polygon 1 & C30/37 & – & $4$ & $480000$ & $400.0$ & $300.0$ \\
Aussparung 2 & – & Polygon 1 & $4$ & $-150000$ & $400.0$ & $300.0$ \\
\hline
\end{longtable}
\noindent\textbf{Fläche des Bruttoquerschnitts}
\begin{equation*}
A = 330000\,\mathrm{mm}^{2}
\end{equation*}
\noindent\textbf{Schwerpunkt, y}
\begin{equation*}
y_S = 400\,\mathrm{mm}
\end{equation*}
\noindent\textbf{Schwerpunkt, z}
\begin{equation*}
z_S = 300\,\mathrm{mm}
\end{equation*}
\noindent\textbf{Trägheitsmoment um y}
\begin{equation*}
I_y = 13275\,\cdot 10^{6}\,\mathrm{mm}^{4}
\end{equation*}
\noindent\textbf{Trägheitsmoment um z}
\begin{equation*}
I_z = 22475\,\cdot 10^{6}\,\mathrm{mm}^{4}
\end{equation*}
\noindent\textbf{Deviationsmoment}
\begin{equation*}
I_{yz} = 0\,\cdot 10^{6}\,\mathrm{mm}^{4}
\end{equation*}
\noindent\textbf{Teilung einer Stablinie: gleiche Felder, auf die Länge gerundet}
\begin{equation*}
s = \frac{L}{m} \qquad m = \mathrm{runde}\left(\frac{L}{s_{gewählt}}\right)
\end{equation*}
\noindent\textbf{Bewehrung}
\begin{longtable}{lllrr}
\hline
Element & $\text{Lage}\ [\mathrm{mm}]$ & Bewehrung & Stahl & $A_s\ [\mathrm{mm}^2]$ \\
\hline
\endhead
Stab 1 & $(760;\ 300)$ & $\varnothing 20$ & B500B & $314$ \\
Linie 1 & $(50;\ 50) \rightarrow (750;\ 50)$ & $6 \varnothing 16\ @\ 140.0\ (\text{gewählt } 150)$ & B500B & $1206$ \\
Linie 2 & $(60;\ 550) \rightarrow (740;\ 550)$ & $4 \varnothing 12\ @\ 136.0$ & B500B & $452$ \\
Linie 3 & $(40;\ 150) \rightarrow (40;\ 450)$ & verschmiert & B500B & $400$ \\
\hline
\end{longtable}
\noindent\textbf{Stahlfläche der Bewehrung}
\begin{equation*}
A_{s,tot} = 2373\,\mathrm{mm}^{2}
\end{equation*}
\noindent\textbf{Bügelquerschnitt je Länge einer Schubwand (n Schnitte)}
\begin{equation*}
\frac{A_{sw}}{s} = \frac{n \cdot \pi\, \varnothing^2 / 4}{s}
\end{equation*}
\noindent\textbf{Schubwände}
\begin{longtable}{llrrlr}
\hline
Wand & $\text{Achse}\ [\mathrm{mm}]$ & $l\ [\mathrm{mm}]$ & $b_w\ [\mathrm{mm}]$ & Bügel & $A_{sw}/s\ [\mathrm{mm}^2/\mathrm{m}]$ \\
\hline
\endhead
Wand 1 & $(75;\ 75) \rightarrow (725;\ 75)$ & $650$ & $150$ & $1 \times \varnothing 10\ @\ 150$ & $524$ \\
Wand 2 & $(725;\ 75) \rightarrow (725;\ 525)$ & $450$ & $150$ & $1 \times \varnothing 10\ @\ 150$ & $524$ \\
Wand 3 & $(725;\ 525) \rightarrow (75;\ 525)$ & $650$ & $150$ & $1 \times \varnothing 10\ @\ 150$ & $524$ \\
Wand 4 & $(75;\ 525) \rightarrow (75;\ 75)$ & $450$ & $150$ & $1 \times \varnothing 10\ @\ 150$ & $524$ \\
\hline
\end{longtable}
\noindent\textbf{Querkraft über die Federn der Wände}
\begin{equation*}
k_i = b_{w,i} \cdot l_i \qquad \sum_i k_i\, e_i e_i^T\, u = V \qquad V_i = k_i\, e_i \cdot u \qquad \Delta T = \sum_i V_i\, r_i
\end{equation*}
\noindent\textbf{Torsion in den Zellen (Bredt, mehrzellig)}
\begin{equation*}
T - \Delta T = \sum_k 2 A_k\, q_k \qquad \oint_k \frac{q}{b_w}\, \mathrm{d}s = 2 A_k\, \theta
\end{equation*}
\noindent\textbf{Widerstand je Länge einer Wand (Fachwerk)} \hfill \normref{SIA 262:2025, 4.3.3.4}
\begin{equation*}
v_{Rd,s} = \frac{A_{sw}}{s}\, f_{sd} \cot\alpha \qquad v_{Rd,c} = b_w\, k_c\, f_{cd} \sin\alpha \cos\alpha \qquad v_{Rd} = \max_\alpha \min(v_{Rd,s};\ v_{Rd,c})
\end{equation*}
\noindent\textbf{Druckfeld}
\begin{longtable}{lr}
\hline
Grenze & Wert \\
\hline
\endhead
$\alpha_{min}$ & $30{}^{\circ}$ \\
$\alpha_{max}$ & $45{}^{\circ}$ \\
$k_c$ & $0.55$ \\
\hline
\end{longtable}
\noindent\textbf{Schubwände und ihre Steifigkeit}
\begin{longtable}{lrrrrl}
\hline
Wand & $l\ [\mathrm{mm}]$ & $b_w\ [\mathrm{mm}]$ & $k = b_w l\ [\mathrm{mm}^2]$ & $A_{sw}/s\ [\mathrm{mm^2/m}]$ & Zelle \\
\hline
\endhead
Wand 1 & $650$ & $150$ & $97500$ & $524$ & 1 \\
Wand 2 & $450$ & $150$ & $67500$ & $524$ & 1 \\
Wand 3 & $650$ & $150$ & $97500$ & $524$ & 1 \\
Wand 4 & $450$ & $150$ & $67500$ & $524$ & 1 \\
\hline
\end{longtable}
\subsubsection{Feld}
\noindent\textbf{Einwirkung}
\begin{longtable}{lrl}
\hline
Einwirkung & Wert &  \\
\hline
\endhead
V\_y,Ed & $0.0$ & kN \\
V\_z,Ed & $250.0$ & kN \\
T\_Ed & $60.0$ & kNm \\
\hline
\end{longtable}
\noindent\textbf{Versatz: Moment der Querkraftanteile um den Schwerpunkt}
\begin{equation*}
\Delta T = \sum_i V_i \cdot r_i = 0\,\mathrm{kNm}
\end{equation*}
\noindent\textbf{Was die Zellen tragen}
\begin{equation*}
T_{Zellen} = T_{Ed} - \Delta T = 60\,\mathrm{kNm} - \left(0\,\mathrm{kNm}\right) = 60\,\mathrm{kNm}
\end{equation*}
\noindent\textbf{Umlauf-Schubfluss der Zellen (Bredt)}
\begin{longtable}{lrr}
\hline
Zelle & $A_k\ [\mathrm{mm}^2]$ & $q_k\ [\mathrm{kN/m}]$ \\
\hline
\endhead
Zelle 1 & $292500$ & $102.6$ \\
\hline
\end{longtable}
\noindent\textbf{Je Wand: Kraft aus der Querkraft, stärkster Schubfluss, Widerstand}
\begin{longtable}{lrrrrrrr}
\hline
Wand & $V_i\ [\mathrm{kN}]$ & $|q_{Ed}|\ [\mathrm{kN/m}]$ & $\alpha\ [{}^{\circ}]$ & $v_{Rd,s}\ [\mathrm{kN/m}]$ & $v_{Rd,c}\ [\mathrm{kN/m}]$ & $v_{Rd}\ [\mathrm{kN/m}]$ & $\alpha_{eff}$ \\
\hline
\endhead
Wand 1 & $0.0$ & $102.6$ & $30$ & $394.3$ & $714.5$ & $394.3$ & $3.84$ \\
Wand 2 & $125.0$ & $380.3$ & $30$ & $394.3$ & $714.5$ & $394.3$ & $1.04$ \\
Wand 3 & $0.0$ & $102.6$ & $30$ & $394.3$ & $714.5$ & $394.3$ & $3.84$ \\
Wand 4 & $-125.0$ & $175.2$ & $30$ & $394.3$ & $714.5$ & $394.3$ & $2.25$ \\
\hline
\end{longtable}
Massgebend: Wand 2.

\noindent\textbf{Erfüllungsgrad}
\begin{equation*}
\alpha_{eff,V,\text{Feld}} = \frac{v_{Rd}}{q_{Ed}} = \frac{394.3\,\mathrm{kN}/\mathrm{m}}{380.3\,\mathrm{kN}/\mathrm{m}} = 1.04 \quad \Rightarrow \quad \text{erfüllt}
\end{equation*}
\noindent\textbf{Längszugkraft aus dem Druckfeld der Wände, je Stück in seiner Mitte}
\begin{longtable}{lrl}
\hline
Längszugkraft & Wert &  \\
\hline
\endhead
$\Delta N = \sum |q|\, l \cot\alpha$ & $664.0$ & kN \\
$\Delta M_y = -\sum \Delta N_i (z_i - z_S)$ & $0.0$ & kNm \\
$\Delta M_z = -\sum \Delta N_i (y_i - y_S)$ & $-52.0$ & kNm \\
\hline
\end{longtable}
\subsubsection{Schief}
\noindent\textbf{Einwirkung}
\begin{longtable}{lrl}
\hline
Einwirkung & Wert &  \\
\hline
\endhead
V\_y,Ed & $100.0$ & kN \\
V\_z,Ed & $0.0$ & kN \\
T\_Ed & $0.0$ & kNm \\
\hline
\end{longtable}
\noindent\textbf{Versatz: Moment der Querkraftanteile um den Schwerpunkt}
\begin{equation*}
\Delta T = \sum_i V_i \cdot r_i = 0\,\mathrm{kNm}
\end{equation*}
\noindent\textbf{Was die Zellen tragen}
\begin{equation*}
T_{Zellen} = T_{Ed} - \Delta T = 0\,\mathrm{kNm} - \left(0\,\mathrm{kNm}\right) = 0\,\mathrm{kNm}
\end{equation*}
\noindent\textbf{Umlauf-Schubfluss der Zellen (Bredt)}
\begin{longtable}{lrr}
\hline
Zelle & $A_k\ [\mathrm{mm}^2]$ & $q_k\ [\mathrm{kN/m}]$ \\
\hline
\endhead
Zelle 1 & $292500$ & $0.0$ \\
\hline
\end{longtable}
\noindent\textbf{Je Wand: Kraft aus der Querkraft, stärkster Schubfluss, Widerstand}
\begin{longtable}{lrrrrrrr}
\hline
Wand & $V_i\ [\mathrm{kN}]$ & $|q_{Ed}|\ [\mathrm{kN/m}]$ & $\alpha\ [{}^{\circ}]$ & $v_{Rd,s}\ [\mathrm{kN/m}]$ & $v_{Rd,c}\ [\mathrm{kN/m}]$ & $v_{Rd}\ [\mathrm{kN/m}]$ & $\alpha_{eff}$ \\
\hline
\endhead
Wand 1 & $50.0$ & $76.9$ & $30$ & $394.3$ & $714.5$ & $394.3$ & $5.13$ \\
Wand 2 & $0.0$ & $0.0$ & $30$ & $394.3$ & $714.5$ & $394.3$ & $\infty$ \\
Wand 3 & $-50.0$ & $76.9$ & $30$ & $394.3$ & $714.5$ & $394.3$ & $5.13$ \\
Wand 4 & $0.0$ & $0.0$ & $30$ & $394.3$ & $714.5$ & $394.3$ & $\infty$ \\
\hline
\end{longtable}
Massgebend: Wand 1.

\noindent\textbf{Erfüllungsgrad}
\begin{equation*}
\alpha_{eff,V,\text{Schief}} = \frac{v_{Rd}}{q_{Ed}} = \frac{394.3\,\mathrm{kN}/\mathrm{m}}{76.9\,\mathrm{kN}/\mathrm{m}} = 5.13 \quad \Rightarrow \quad \text{erfüllt}
\end{equation*}
\noindent\textbf{Längszugkraft aus dem Druckfeld der Wände, je Stück in seiner Mitte}
\begin{longtable}{lrl}
\hline
Längszugkraft & Wert &  \\
\hline
\endhead
$\Delta N = \sum |q|\, l \cot\alpha$ & $173.2$ & kN \\
$\Delta M_y = -\sum \Delta N_i (z_i - z_S)$ & $0.0$ & kNm \\
$\Delta M_z = -\sum \Delta N_i (y_i - y_S)$ & $0.0$ & kNm \\
\hline
\end{longtable}
\noindent\textbf{Dehnungsebene und innere Kräfte}
\begin{equation*}
\varepsilon(y, z) = \varepsilon_0 + \kappa \cdot v \qquad v = n_y (y - y_S) + n_z (z - z_S) \qquad N = \int_A \sigma\,\mathrm{d}A \qquad M_y = -\int_A \sigma\,(z - z_S)\,\mathrm{d}A \qquad M_z = -\int_A \sigma\,(y - y_S)\,\mathrm{d}A
\end{equation*}
\noindent\textbf{Erfüllungsgrad bei fester Normalkraft und fester Richtung}
\begin{equation*}
\alpha_{eff} = \frac{M_{Rd}}{M_{Ed}} \qquad M = \sqrt{M_y^2 + M_z^2}
\end{equation*}
\noindent\textbf{C30/37: Spannungsblock 0.85·x, Bemessungswerte}
\begin{equation*}
\sigma_c(\varepsilon_c) = \begin{cases} 0 & \varepsilon_c > -(1 - 0.85) \cdot \varepsilon_{c2d} \\[1ex] -f_{cd} & -\varepsilon_{c2d} \le \varepsilon_c \le -(1 - 0.85) \cdot \varepsilon_{c2d} \end{cases}
\end{equation*}
\noindent\textbf{Werte C30/37}
\begin{longtable}{lr}
\hline
Kennwert & Wert \\
\hline
\endhead
$f_{cd}$ & $20\,\mathrm{N}/\mathrm{mm}^{2}$ \\
$\varepsilon_{c1d}$ & $2\,\text{‰}$ \\
$\varepsilon_{c2d}$ & $3.5\,\text{‰}$ \\
\hline
\end{longtable}
\noindent\textbf{B500B: bilineare Beziehung, Bemessungswerte} \hfill \normref{SIA 262:2025, 4.2.2.4}
\begin{equation*}
\sigma_s(\varepsilon_s) = \begin{cases} \min(E_s\,\varepsilon_s;\ f_{yd}) & \varepsilon_s \ge 0 \\[1ex] \max(E_s\,\varepsilon_s;\ -f_{yd}^{-}) & \varepsilon_s < 0 \end{cases} \qquad |\varepsilon_s| \le \varepsilon_{ud}
\end{equation*}
\noindent\textbf{Werte B500B}
\begin{longtable}{lr}
\hline
Kennwert & Wert \\
\hline
\endhead
$f_{yd}$ & $435\,\mathrm{N}/\mathrm{mm}^{2}$ \\
$f_{yd}^{-}$ & $435\,\mathrm{N}/\mathrm{mm}^{2}$ \\
$E_s$ & $200000\,\mathrm{N}/\mathrm{mm}^{2}$ \\
$\varepsilon_{ud}$ & $4.5\,\%$ \\
\hline
\end{longtable}
\noindent\textbf{Reine Normalkraft}
\begin{longtable}{lr}
\hline
Grenze & Wert \\
\hline
\endhead
reiner Zug, gleichmässig & $N_{Rd}^{+} = 1031.7\,\mathrm{kN}$ \\
reiner Druck, gleichmässig & $N_{Rd}^{-} = -7501.7\,\mathrm{kN}$ \\
\hline
\end{longtable}
\subsubsection{Feld}
\noindent\textbf{Einwirkung}
\begin{longtable}{lrl}
\hline
Einwirkung & Wert &  \\
\hline
\endhead
N\_Ed & $0.0$ & kN \\
M\_y,Ed & $300.0$ & kNm \\
M\_z,Ed & $0.0$ & kNm \\
ΔN aus Querkraft und Torsion & $664.0$ & kN \\
ΔM\_y & $0.0$ & kNm \\
ΔM\_z & $-52.0$ & kNm \\
N (nachgewiesen) & $664.0$ & kN \\
M\_y (nachgewiesen) & $300.0$ & kNm \\
M\_z (nachgewiesen) & $-52.0$ & kNm \\
\hline
\end{longtable}
Bruchzustand: Zugrichtung n = (0.046; -0.999), -87.3° gegen die y-Achse. Die Nulllinie steht senkrecht zur Zugrichtung, 272.7 mm auf der Druckseite des Schwerpunkts.

\noindent\textbf{Dehnungsebene}
\begin{longtable}{rr}
\hline
$\varepsilon_0\ [\text{‰}]$ & $\kappa\ [\mathrm{1/m}]$ \\
\hline
\endhead
$20.965$ & $0.07689$ \\
\hline
\end{longtable}
\noindent\textbf{Kräfte im Bruchzustand -- Beton als Resultierende, Bewehrung je Element}
\begin{longtable}{lrrrrr}
\hline
Teil & $F\ [\mathrm{kN}]$ & $y\ [\mathrm{mm}]$ & $z\ [\mathrm{mm}]$ & $\varepsilon\ [\text{‰}]$ & $\sigma\ [\mathrm{N/mm^2}]$ \\
\hline
\endhead
C30/37 & $-323.4$ & $277.4$ & $587.0$ & – & – \\
Stab 1 & $136.6$ & $760.0$ & $300.0$ & $22.25$ & $435$ \\
Linie 1 & $524.5$ & $400.0$ & $50.0$ & $38.92 \ldots 41.41$ & $435$ \\
Linie 2 & $152.4$ & $439.3$ & $550.0$ & $1.04 \ldots 2.49$ & $207 \ldots 435$ \\
Linie 3 & $173.9$ & $40.0$ & $300.0$ & $8.88 \ldots 30.49$ & $435$ \\
\hline
\end{longtable}
\noindent\textbf{Probe: Normalkraft}
\begin{equation*}
N_{Rd} = \sum F_i = 664\,\mathrm{kN}
\end{equation*}
\noindent\textbf{Probe: Moment um y}
\begin{equation*}
M_{y,Rd} = -\sum F_i \cdot (z_i - z_S) = 185.9\,\mathrm{kNm}
\end{equation*}
\noindent\textbf{Probe: Moment um z}
\begin{equation*}
M_{z,Rd} = -\sum F_i \cdot (y_i - y_S) = -32.2\,\mathrm{kNm}
\end{equation*}
\noindent\textbf{Widerstand in der Richtung der Einwirkung}
\begin{equation*}
M_{Rd} = \sqrt{M_{y,Rd}^{2} + M_{z,Rd}^{2}} = \sqrt{\left(185.9\,\mathrm{kNm}\right)^{2} + \left(-32.2\,\mathrm{kNm}\right)^{2}} = 188.6\,\mathrm{kNm}
\end{equation*}
\noindent\textbf{Erfüllungsgrad}
\begin{equation*}
\alpha_{eff,\text{Feld}} = \frac{M_{Rd}}{M_{Ed}} = \frac{188.6\,\mathrm{kNm}}{304.5\,\mathrm{kNm}} = 0.62 \quad \Rightarrow \quad \text{NICHT erfüllt}
\end{equation*}
\subsubsection{Schief}
\noindent\textbf{Einwirkung}
\begin{longtable}{lrl}
\hline
Einwirkung & Wert &  \\
\hline
\endhead
N\_Ed & $-1500.0$ & kN \\
M\_y,Ed & $250.0$ & kNm \\
M\_z,Ed & $150.0$ & kNm \\
ΔN aus Querkraft und Torsion & $173.2$ & kN \\
ΔM\_y & $0.0$ & kNm \\
ΔM\_z & $0.0$ & kNm \\
N (nachgewiesen) & $-1326.8$ & kN \\
M\_y (nachgewiesen) & $250.0$ & kNm \\
M\_z (nachgewiesen) & $150.0$ & kNm \\
\hline
\end{longtable}
Bruchzustand: Zugrichtung n = (-0.333; -0.943), -109.4° gegen die y-Achse. Die Nulllinie steht senkrecht zur Zugrichtung, 120.3 mm auf der Druckseite des Schwerpunkts.

\noindent\textbf{Dehnungsebene}
\begin{longtable}{rr}
\hline
$\varepsilon_0\ [\text{‰}]$ & $\kappa\ [\mathrm{1/m}]$ \\
\hline
\endhead
$1.424$ & $0.01184$ \\
\hline
\end{longtable}
\noindent\textbf{Kräfte im Bruchzustand -- Beton als Resultierende, Bewehrung je Element}
\begin{longtable}{lrrrrr}
\hline
Teil & $F\ [\mathrm{kN}]$ & $y\ [\mathrm{mm}]$ & $z\ [\mathrm{mm}]$ & $\varepsilon\ [\text{‰}]$ & $\sigma\ [\mathrm{N/mm^2}]$ \\
\hline
\endhead
C30/37 & $-1898.5$ & $545.7$ & $516.1$ & – & – \\
Stab 1 & $0.4$ & $760.0$ & $300.0$ & $0.01$ & $1$ \\
Linie 1 & $524.5$ & $400.0$ & $50.0$ & $2.84 \ldots 5.59$ & $435$ \\
Linie 2 & $-115.0$ & $470.8$ & $550.0$ & $-2.17 \ldots -0.56$ & $-434 \ldots -113$ \\
Linie 3 & $161.8$ & $40.0$ & $291.1$ & $1.27 \ldots 4.41$ & $255 \ldots 435$ \\
\hline
\end{longtable}
\noindent\textbf{Probe: Normalkraft}
\begin{equation*}
N_{Rd} = \sum F_i = -1326.8\,\mathrm{kN}
\end{equation*}
\noindent\textbf{Probe: Moment um y}
\begin{equation*}
M_{y,Rd} = -\sum F_i \cdot (z_i - z_S) = 571.6\,\mathrm{kNm}
\end{equation*}
\noindent\textbf{Probe: Moment um z}
\begin{equation*}
M_{z,Rd} = -\sum F_i \cdot (y_i - y_S) = 343\,\mathrm{kNm}
\end{equation*}
\noindent\textbf{Widerstand in der Richtung der Einwirkung}
\begin{equation*}
M_{Rd} = \sqrt{M_{y,Rd}^{2} + M_{z,Rd}^{2}} = \sqrt{\left(571.6\,\mathrm{kNm}\right)^{2} + \left(343\,\mathrm{kNm}\right)^{2}} = 666.6\,\mathrm{kNm}
\end{equation*}
\noindent\textbf{Erfüllungsgrad}
\begin{equation*}
\alpha_{eff,\text{Schief}} = \frac{M_{Rd}}{M_{Ed}} = \frac{666.6\,\mathrm{kNm}}{291.5\,\mathrm{kNm}} = 2.29 \quad \Rightarrow \quad \text{erfüllt}
\end{equation*}
\noindent\textbf{Duktilität} \hfill \normref{SIA 262:2025, 4.1.4.2.5}
\begin{equation*}
\frac{x}{d} \le \left(\frac{x}{d}\right)_{max}
\end{equation*}
\subsubsection{Zug unten}
\noindent\textbf{Bezogene Druckzonenhöhe bei reiner Biegung}
\begin{equation*}
x/d = \frac{x}{d} = \frac{64.8\,\mathrm{mm}}{560.3\,\mathrm{mm}} = 0.116 \quad \le \quad 0.35 \quad \Rightarrow \quad \text{erfüllt}
\end{equation*}
\noindent\textbf{Erfüllungsgrad}
\begin{equation*}
\alpha_{eff,x/d,\text{Zug unten}} = \frac{\left(x/d\right)_{max}}{x/d} = \frac{0.35}{0.116} = 3.02
\end{equation*}
\subsubsection{Zug oben}
\noindent\textbf{Bezogene Druckzonenhöhe bei reiner Biegung}
\begin{equation*}
x/d = \frac{x}{d} = \frac{47.4\,\mathrm{mm}}{554.1\,\mathrm{mm}} = 0.086 \quad \le \quad 0.35 \quad \Rightarrow \quad \text{erfüllt}
\end{equation*}
\noindent\textbf{Erfüllungsgrad}
\begin{equation*}
\alpha_{eff,x/d,\text{Zug oben}} = \frac{\left(x/d\right)_{max}}{x/d} = \frac{0.35}{0.086} = 4.09
\end{equation*}
\subsubsection{Zug links}
\noindent\textbf{Bezogene Druckzonenhöhe bei reiner Biegung}
\begin{equation*}
x/d = \frac{x}{d} = \frac{109\,\mathrm{mm}}{807.7\,\mathrm{mm}} = 0.135 \quad \le \quad 0.35 \quad \Rightarrow \quad \text{erfüllt}
\end{equation*}
\noindent\textbf{Erfüllungsgrad}
\begin{equation*}
\alpha_{eff,x/d,\text{Zug links}} = \frac{\left(x/d\right)_{max}}{x/d} = \frac{0.35}{0.135} = 2.59
\end{equation*}
\subsubsection{Zug rechts}
\noindent\textbf{Bezogene Druckzonenhöhe bei reiner Biegung}
\begin{equation*}
x/d = \frac{x}{d} = \frac{103.5\,\mathrm{mm}}{789.9\,\mathrm{mm}} = 0.131 \quad \le \quad 0.35 \quad \Rightarrow \quad \text{erfüllt}
\end{equation*}
\noindent\textbf{Erfüllungsgrad}
\begin{equation*}
\alpha_{eff,x/d,\text{Zug rechts}} = \frac{\left(x/d\right)_{max}}{x/d} = \frac{0.35}{0.131} = 2.67
\end{equation*}
Massgebend: Zug links (kleinster Erfüllungsgrad).

\noindent\textbf{Spannung im ungerissenen Querschnitt}
\begin{equation*}
\sigma(y, z) = -\frac{M_y\,(I_z z' - I_{yz} y') + M_z\,(I_y y' - I_{yz} z')}{I_y I_z - I_{yz}^2}\qquad y' = y - y_S,\ z' = z - z_S
\end{equation*}
\noindent\textbf{Sprödes Versagen} \hfill \normref{SIA 262:2025, 4.4.1.3}
\begin{equation*}
M_{Rd}(N = 0) \ge M_{Riss} = f_{ct,eff} \cdot W
\end{equation*}
\subsubsection{Zug unten}
Gezogenster Betonpunkt: y = 0 mm, z = 0 mm (C30/37).

\noindent\textbf{Beiwert für die Höhe in Biegerichtung} \hfill \normref{SIA 262:2025, 4.4.1.3}
\begin{equation*}
k_t = \frac{1}{1 + 0.5 \cdot h/3} = \frac{1}{1 + 0.5 \cdot 0.6/3} = 0.909 \quad \left(h\ \text{in}\ \mathrm{m}\right)
\end{equation*}
\noindent\textbf{Wirksame Zugfestigkeit}
\begin{equation*}
f_{ct,eff} = k_t \cdot f_{ctm} = 0.909 \cdot 2.9\,\mathrm{N}/\mathrm{mm}^{2} = 2.64\,\mathrm{N}/\mathrm{mm}^{2}
\end{equation*}
\noindent\textbf{Widerstandsmoment zum gezogensten Betonpunkt}
\begin{equation*}
\begin{aligned}
  W &= \frac{I_y \cdot I_z - I_{yz}^{2}}{I_{yz} \cdot y' - I_z \cdot z'} \\
  &= \frac{13275\,\cdot 10^{6}\,\mathrm{mm}^{4} \cdot 22475\,\cdot 10^{6}\,\mathrm{mm}^{4} - \left(0\,\cdot 10^{6}\,\mathrm{mm}^{4}\right)^{2}}{0\,\cdot 10^{6}\,\mathrm{mm}^{4} \cdot \left(-400\,\mathrm{mm}\right) - 22475\,\cdot 10^{6}\,\mathrm{mm}^{4} \cdot \left(-300\,\mathrm{mm}\right)} \\
  &= 44.25\,\cdot 10^{6}\,\mathrm{mm}^{3}
\end{aligned}
\end{equation*}
\noindent\textbf{Rissmoment des ungerissenen Querschnitts}
\begin{equation*}
M_{Riss} = f_{ct,eff} \cdot W = 2.64\,\mathrm{N}/\mathrm{mm}^{2} \cdot 44.25\,\cdot 10^{6}\,\mathrm{mm}^{3} = 116.7\,\mathrm{kNm}
\end{equation*}
\noindent\textbf{Biegewiderstand bei N = 0 in dieser Richtung}
\begin{equation*}
M_{Rd} = 359.8\,\mathrm{kNm}
\end{equation*}
\noindent\textbf{Erfüllungsgrad}
\begin{equation*}
\alpha_{eff,Riss,\text{Zug unten}} = \frac{M_{Rd}}{M_{Riss}} = \frac{359.8\,\mathrm{kNm}}{116.7\,\mathrm{kNm}} = 3.08
\end{equation*}
\subsubsection{Zug oben}
Gezogenster Betonpunkt: y = 800 mm, z = 600 mm (C30/37).

\noindent\textbf{Beiwert für die Höhe in Biegerichtung} \hfill \normref{SIA 262:2025, 4.4.1.3}
\begin{equation*}
k_t = \frac{1}{1 + 0.5 \cdot h/3} = \frac{1}{1 + 0.5 \cdot 0.6/3} = 0.909 \quad \left(h\ \text{in}\ \mathrm{m}\right)
\end{equation*}
\noindent\textbf{Wirksame Zugfestigkeit}
\begin{equation*}
f_{ct,eff} = k_t \cdot f_{ctm} = 0.909 \cdot 2.9\,\mathrm{N}/\mathrm{mm}^{2} = 2.64\,\mathrm{N}/\mathrm{mm}^{2}
\end{equation*}
\noindent\textbf{Widerstandsmoment zum gezogensten Betonpunkt}
\begin{equation*}
\begin{aligned}
  W &= \frac{I_y \cdot I_z - I_{yz}^{2}}{I_z \cdot z' - I_{yz} \cdot y'} \\
  &= \frac{13275\,\cdot 10^{6}\,\mathrm{mm}^{4} \cdot 22475\,\cdot 10^{6}\,\mathrm{mm}^{4} - \left(0\,\cdot 10^{6}\,\mathrm{mm}^{4}\right)^{2}}{22475\,\cdot 10^{6}\,\mathrm{mm}^{4} \cdot 300\,\mathrm{mm} - 0\,\cdot 10^{6}\,\mathrm{mm}^{4} \cdot 400\,\mathrm{mm}} \\
  &= 44.25\,\cdot 10^{6}\,\mathrm{mm}^{3}
\end{aligned}
\end{equation*}
\noindent\textbf{Rissmoment des ungerissenen Querschnitts}
\begin{equation*}
M_{Riss} = f_{ct,eff} \cdot W = 2.64\,\mathrm{N}/\mathrm{mm}^{2} \cdot 44.25\,\cdot 10^{6}\,\mathrm{mm}^{3} = 116.7\,\mathrm{kNm}
\end{equation*}
\noindent\textbf{Biegewiderstand bei N = 0 in dieser Richtung}
\begin{equation*}
M_{Rd} = 194.7\,\mathrm{kNm}
\end{equation*}
\noindent\textbf{Erfüllungsgrad}
\begin{equation*}
\alpha_{eff,Riss,\text{Zug oben}} = \frac{M_{Rd}}{M_{Riss}} = \frac{194.7\,\mathrm{kNm}}{116.7\,\mathrm{kNm}} = 1.67
\end{equation*}
\subsubsection{Zug links}
Gezogenster Betonpunkt: y = 0 mm, z = 0 mm (C30/37).

\noindent\textbf{Beiwert für die Höhe in Biegerichtung} \hfill \normref{SIA 262:2025, 4.4.1.3}
\begin{equation*}
k_t = \frac{1}{1 + 0.5 \cdot h/3} = \frac{1}{1 + 0.5 \cdot 0.8/3} = 0.882 \quad \left(h\ \text{in}\ \mathrm{m}\right)
\end{equation*}
\noindent\textbf{Wirksame Zugfestigkeit}
\begin{equation*}
f_{ct,eff} = k_t \cdot f_{ctm} = 0.882 \cdot 2.9\,\mathrm{N}/\mathrm{mm}^{2} = 2.56\,\mathrm{N}/\mathrm{mm}^{2}
\end{equation*}
\noindent\textbf{Widerstandsmoment zum gezogensten Betonpunkt}
\begin{equation*}
\begin{aligned}
  W &= \frac{I_y \cdot I_z - I_{yz}^{2}}{I_{yz} \cdot z' - I_y \cdot y'} \\
  &= \frac{13275\,\cdot 10^{6}\,\mathrm{mm}^{4} \cdot 22475\,\cdot 10^{6}\,\mathrm{mm}^{4} - \left(0\,\cdot 10^{6}\,\mathrm{mm}^{4}\right)^{2}}{0\,\cdot 10^{6}\,\mathrm{mm}^{4} \cdot \left(-300\,\mathrm{mm}\right) - 13275\,\cdot 10^{6}\,\mathrm{mm}^{4} \cdot \left(-400\,\mathrm{mm}\right)} \\
  &= 56.19\,\cdot 10^{6}\,\mathrm{mm}^{3}
\end{aligned}
\end{equation*}
\noindent\textbf{Rissmoment des ungerissenen Querschnitts}
\begin{equation*}
M_{Riss} = f_{ct,eff} \cdot W = 2.56\,\mathrm{N}/\mathrm{mm}^{2} \cdot 56.19\,\cdot 10^{6}\,\mathrm{mm}^{3} = 143.8\,\mathrm{kNm}
\end{equation*}
\noindent\textbf{Biegewiderstand bei N = 0 in dieser Richtung}
\begin{equation*}
M_{Rd} = 388.5\,\mathrm{kNm}
\end{equation*}
\noindent\textbf{Erfüllungsgrad}
\begin{equation*}
\alpha_{eff,Riss,\text{Zug links}} = \frac{M_{Rd}}{M_{Riss}} = \frac{388.5\,\mathrm{kNm}}{143.8\,\mathrm{kNm}} = 2.70
\end{equation*}
\subsubsection{Zug rechts}
Gezogenster Betonpunkt: y = 800 mm, z = 0 mm (C30/37).

\noindent\textbf{Beiwert für die Höhe in Biegerichtung} \hfill \normref{SIA 262:2025, 4.4.1.3}
\begin{equation*}
k_t = \frac{1}{1 + 0.5 \cdot h/3} = \frac{1}{1 + 0.5 \cdot 0.8/3} = 0.882 \quad \left(h\ \text{in}\ \mathrm{m}\right)
\end{equation*}
\noindent\textbf{Wirksame Zugfestigkeit}
\begin{equation*}
f_{ct,eff} = k_t \cdot f_{ctm} = 0.882 \cdot 2.9\,\mathrm{N}/\mathrm{mm}^{2} = 2.56\,\mathrm{N}/\mathrm{mm}^{2}
\end{equation*}
\noindent\textbf{Widerstandsmoment zum gezogensten Betonpunkt}
\begin{equation*}
\begin{aligned}
  W &= \frac{I_y \cdot I_z - I_{yz}^{2}}{I_y \cdot y' - I_{yz} \cdot z'} \\
  &= \frac{13275\,\cdot 10^{6}\,\mathrm{mm}^{4} \cdot 22475\,\cdot 10^{6}\,\mathrm{mm}^{4} - \left(0\,\cdot 10^{6}\,\mathrm{mm}^{4}\right)^{2}}{13275\,\cdot 10^{6}\,\mathrm{mm}^{4} \cdot 400\,\mathrm{mm} - 0\,\cdot 10^{6}\,\mathrm{mm}^{4} \cdot \left(-300\,\mathrm{mm}\right)} \\
  &= 56.19\,\cdot 10^{6}\,\mathrm{mm}^{3}
\end{aligned}
\end{equation*}
\noindent\textbf{Rissmoment des ungerissenen Querschnitts}
\begin{equation*}
M_{Riss} = f_{ct,eff} \cdot W = 2.56\,\mathrm{N}/\mathrm{mm}^{2} \cdot 56.19\,\cdot 10^{6}\,\mathrm{mm}^{3} = 143.8\,\mathrm{kNm}
\end{equation*}
\noindent\textbf{Biegewiderstand bei N = 0 in dieser Richtung}
\begin{equation*}
M_{Rd} = 362.5\,\mathrm{kNm}
\end{equation*}
\noindent\textbf{Erfüllungsgrad}
\begin{equation*}
\alpha_{eff,Riss,\text{Zug rechts}} = \frac{M_{Rd}}{M_{Riss}} = \frac{362.5\,\mathrm{kNm}}{143.8\,\mathrm{kNm}} = 2.52
\end{equation*}
Massgebend: Zug oben (kleinster Erfüllungsgrad).

\subsection{Querschnittsanalyse: Stütze}
\noindent\textbf{Fläche und Schwerpunkt eines Polygons (Satz von Gauss)}
\begin{equation*}
A = \frac{1}{2} \sum_i c_i \qquad y_S = \frac{1}{6A} \sum_i (y_i + y_{i+1})\, c_i \qquad z_S = \frac{1}{6A} \sum_i (z_i + z_{i+1})\, c_i \qquad c_i = y_i\, z_{i+1} - y_{i+1}\, z_i
\end{equation*}
\noindent\textbf{Trägheitsmomente eines Polygons, bezogen auf den Schwerpunkt}
\begin{equation*}
I_y = \frac{1}{12} \sum_i (z_i^2 + z_i z_{i+1} + z_{i+1}^2)\, c_i - A\, z_S^2\qquad I_z = \frac{1}{12} \sum_i (y_i^2 + y_i y_{i+1} + y_{i+1}^2)\, c_i - A\, y_S^2
\end{equation*}
\noindent\textbf{Polygone des Querschnitts}
\begin{longtable}{lllrrrr}
\hline
Polygon & Material & liegt in & Ecken & $\pm A\ [\mathrm{mm}^2]$ & $y_S\ [\mathrm{mm}]$ & $z_S\ [\mathrm{mm}]$ \\
\hline
\endhead
Polygon 1 & C30/37 & – & $4$ & $160000$ & $200.0$ & $200.0$ \\
\hline
\end{longtable}
\noindent\textbf{Fläche des Bruttoquerschnitts}
\begin{equation*}
A = 160000\,\mathrm{mm}^{2}
\end{equation*}
\noindent\textbf{Schwerpunkt, y}
\begin{equation*}
y_S = 200\,\mathrm{mm}
\end{equation*}
\noindent\textbf{Schwerpunkt, z}
\begin{equation*}
z_S = 200\,\mathrm{mm}
\end{equation*}
\noindent\textbf{Trägheitsmoment um y}
\begin{equation*}
I_y = 2133.3\,\cdot 10^{6}\,\mathrm{mm}^{4}
\end{equation*}
\noindent\textbf{Trägheitsmoment um z}
\begin{equation*}
I_z = 2133.3\,\cdot 10^{6}\,\mathrm{mm}^{4}
\end{equation*}
\noindent\textbf{Deviationsmoment}
\begin{equation*}
I_{yz} = 0\,\cdot 10^{6}\,\mathrm{mm}^{4}
\end{equation*}
\noindent\textbf{Teilung einer Stablinie: gleiche Felder, auf die Länge gerundet}
\begin{equation*}
s = \frac{L}{m} \qquad m = \mathrm{runde}\left(\frac{L}{s_{gewählt}}\right)
\end{equation*}
\noindent\textbf{Bewehrung}
\begin{longtable}{lllrr}
\hline
Element & $\text{Lage}\ [\mathrm{mm}]$ & Bewehrung & Stahl & $A_s\ [\mathrm{mm}^2]$ \\
\hline
\endhead
Linie 1 & $(50;\ 50) \rightarrow (350;\ 50)$ & $3 \varnothing 20\ @\ 150.0$ & B500B & $942$ \\
Linie 2 & $(50;\ 350) \rightarrow (350;\ 350)$ & $3 \varnothing 20\ @\ 150.0$ & B500B & $942$ \\
\hline
\end{longtable}
\noindent\textbf{Stahlfläche der Bewehrung}
\begin{equation*}
A_{s,tot} = 1885\,\mathrm{mm}^{2}
\end{equation*}
\noindent\textbf{Dehnungsebene und innere Kräfte}
\begin{equation*}
\varepsilon(y, z) = \varepsilon_0 + \kappa \cdot v \qquad v = n_y (y - y_S) + n_z (z - z_S) \qquad N = \int_A \sigma\,\mathrm{d}A \qquad M_y = -\int_A \sigma\,(z - z_S)\,\mathrm{d}A \qquad M_z = -\int_A \sigma\,(y - y_S)\,\mathrm{d}A
\end{equation*}
\noindent\textbf{Erfüllungsgrad bei fester Normalkraft und fester Richtung}
\begin{equation*}
\alpha_{eff} = \frac{M_{Rd}}{M_{Ed}} \qquad M = \sqrt{M_y^2 + M_z^2}
\end{equation*}
\noindent\textbf{C30/37: Parabel-Rechteck-Beziehung, charakteristische Werte} \hfill \normref{SIA 262:2025, 4.2.1.6}
\begin{equation*}
\sigma_c(\varepsilon_c) = \begin{cases} -f_{ck} \cdot \dfrac{k_\sigma\,\eta - \eta^2}{1 + (k_\sigma - 2)\,\eta} & 0 \le |\varepsilon_c| \le \varepsilon_{c1d},\ \eta = \dfrac{|\varepsilon_c|}{\varepsilon_{c1d}} \\[2ex] -f_{ck} & \varepsilon_{c1d} < |\varepsilon_c| \le \varepsilon_{c2d} \\[1ex] 0 & \varepsilon_c > 0 \quad (\text{Zug, gerissen}) \end{cases}
\end{equation*}
\noindent\textbf{Krümmungsbeiwert mit charakteristischer Festigkeit} \hfill \normref{SIA 262:2025, 4.2.1.6}
\begin{equation*}
k_{\sigma,k} = \frac{E_{cd}}{400 \cdot f_{ck}} = \frac{33620\,\mathrm{N}/\mathrm{mm}^{2}}{400 \cdot 30\,\mathrm{N}/\mathrm{mm}^{2}} = 2.802
\end{equation*}
\noindent\textbf{Werte C30/37}
\begin{longtable}{lr}
\hline
Kennwert & Wert \\
\hline
\endhead
$f_{ck}$ & $30\,\mathrm{N}/\mathrm{mm}^{2}$ \\
$\varepsilon_{c1d}$ & $2\,\text{‰}$ \\
$\varepsilon_{c2d}$ & $3.5\,\text{‰}$ \\
\hline
\end{longtable}
\noindent\textbf{B500B: bilineare Beziehung, charakteristische Werte} \hfill \normref{SIA 262:2025, 4.2.2.4}
\begin{equation*}
\sigma_s(\varepsilon_s) = \begin{cases} \min(E_s\,\varepsilon_s;\ f_{yk}) & \varepsilon_s \ge 0 \\[1ex] \max(E_s\,\varepsilon_s;\ -f_{yk}^{-}) & \varepsilon_s < 0 \end{cases} \qquad |\varepsilon_s| \le \varepsilon_{ud}
\end{equation*}
\noindent\textbf{Werte B500B}
\begin{longtable}{lr}
\hline
Kennwert & Wert \\
\hline
\endhead
$f_{yk}$ & $500\,\mathrm{N}/\mathrm{mm}^{2}$ \\
$f_{yk}^{-}$ & $500\,\mathrm{N}/\mathrm{mm}^{2}$ \\
$E_s$ & $200000\,\mathrm{N}/\mathrm{mm}^{2}$ \\
$\varepsilon_{ud}$ & $4.5\,\%$ \\
\hline
\end{longtable}
\noindent\textbf{Reine Normalkraft}
\begin{longtable}{lr}
\hline
Grenze & Wert \\
\hline
\endhead
reiner Zug, gleichmässig & $N_{Rd}^{+} = 942.5\,\mathrm{kN}$ \\
reiner Druck, gleichmässig & $N_{Rd}^{-} = -5497.4\,\mathrm{kN}$ \\
\hline
\end{longtable}
\subsubsection{Druck}
\noindent\textbf{Einwirkung}
\begin{longtable}{lrl}
\hline
Einwirkung & Wert &  \\
\hline
\endhead
N\_Ed & $-2500.0$ & kN \\
M\_y,Ed & $120.0$ & kNm \\
\hline
\end{longtable}
Bruchzustand: Zugrichtung n = (0.000; -1.000), -90.0° gegen die y-Achse. Die Nulllinie steht senkrecht zur Zugrichtung, 36.0 mm auf der Zugseite des Schwerpunkts.

\noindent\textbf{Dehnungsebene}
\begin{longtable}{rr}
\hline
$\varepsilon_0\ [\text{‰}]$ & $\kappa\ [\mathrm{1/m}]$ \\
\hline
\endhead
$-0.534$ & $0.01483$ \\
\hline
\end{longtable}
\noindent\textbf{Kräfte im Bruchzustand -- Beton als Resultierende, Bewehrung je Element}
\begin{longtable}{lrrrrr}
\hline
Teil & $F\ [\mathrm{kN}]$ & $y\ [\mathrm{mm}]$ & $z\ [\mathrm{mm}]$ & $\varepsilon\ [\text{‰}]$ & $\sigma\ [\mathrm{N/mm^2}]$ \\
\hline
\endhead
C30/37 & $-2375.7$ & $200.0$ & $298.8$ & – & – \\
Linie 1 & $318.7$ & $200.0$ & $50.0$ & $1.69$ & $338$ \\
Linie 2 & $-443.0$ & $200.0$ & $350.0$ & $-2.76$ & $-500$ \\
\hline
\end{longtable}
\noindent\textbf{Probe: Normalkraft}
\begin{equation*}
N_{Rd} = \sum F_i = -2500\,\mathrm{kN}
\end{equation*}
\noindent\textbf{Probe: Moment um y}
\begin{equation*}
M_{y,Rd} = -\sum F_i \cdot (z_i - z_S) = 348.9\,\mathrm{kNm}
\end{equation*}
\noindent\textbf{Erfüllungsgrad}
\begin{equation*}
\alpha_{eff,\text{Druck}} = \frac{M_{Rd}}{M_{Ed}} = \frac{348.9\,\mathrm{kNm}}{120\,\mathrm{kNm}} = 2.91 \quad \Rightarrow \quad \text{erfüllt}
\end{equation*}
\subsubsection{Zug}
\noindent\textbf{Einwirkung}
\begin{longtable}{lrl}
\hline
Einwirkung & Wert &  \\
\hline
\endhead
N\_Ed & $400.0$ & kN \\
M\_y,Ed & $0.0$ & kNm \\
\hline
\end{longtable}
\noindent\textbf{Erfüllungsgrad}
\begin{equation*}
\alpha_{eff,\text{Zug}} = \frac{N_{Rd}^{+}}{N_{Ed}} = \frac{942.5\,\mathrm{kN}}{400\,\mathrm{kN}} = 2.36 \quad \Rightarrow \quad \text{erfüllt}
\end{equation*}
\section{Verwendete Formeln}
\subsection{Querschnittsanalyse}
Der Querschnitt ist gezeichnet: Polygone, jedes mit seinem Material. Liegt ein Polygon ganz in einem anderen, ersetzt es dieses dort -- eine Aussparung nimmt die Fläche weg. y zeigt nach rechts, z nach oben; alle Masse in Millimetern.

\noindent\textbf{Fläche und Schwerpunkt eines Polygons (Satz von Gauss)}
\begin{equation*}
A = \frac{1}{2} \sum_i c_i \qquad y_S = \frac{1}{6A} \sum_i (y_i + y_{i+1})\, c_i \qquad z_S = \frac{1}{6A} \sum_i (z_i + z_{i+1})\, c_i \qquad c_i = y_i\, z_{i+1} - y_{i+1}\, z_i
\end{equation*}
\noindent\textbf{Trägheitsmomente eines Polygons, bezogen auf den Schwerpunkt}
\begin{equation*}
I_y = \frac{1}{12} \sum_i (z_i^2 + z_i z_{i+1} + z_{i+1}^2)\, c_i - A\, z_S^2\qquad I_z = \frac{1}{12} \sum_i (y_i^2 + y_i y_{i+1} + y_{i+1}^2)\, c_i - A\, y_S^2
\end{equation*}
\noindent\textbf{Teilung einer Stablinie: gleiche Felder, auf die Länge gerundet}
\begin{equation*}
s = \frac{L}{m} \qquad m = \mathrm{runde}\left(\frac{L}{s_{gewählt}}\right)
\end{equation*}
\noindent\textbf{Bügelquerschnitt je Länge einer Schubwand (n Schnitte)}
\begin{equation*}
\frac{A_{sw}}{s} = \frac{n \cdot \pi\, \varnothing^2 / 4}{s}
\end{equation*}
\subsection{Beton}
\noindent\textbf{Beiwert zur Berücksichtigung der Festigkeitsminderung} \hfill \normref{SIA 262:2025, 2.4.2.3}
\begin{equation*}
\eta_{fc} = \min\left[\left(\frac{40\,\mathrm{N}/\mathrm{mm}^{2}}{f_{ck}}\right)^{1/3};\ 1.0\right]
\end{equation*}
\noindent\textbf{Bemessungswert der Betondruckfestigkeit} \hfill \normref{SIA 262:2025, 2.4.2.3}
\begin{equation*}
f_{cd} = \frac{\eta_{fc} \cdot f_{ck}}{\gamma_c}
\end{equation*}
\noindent\textbf{Mittelwert der Zylinderdruckfestigkeit} \hfill \normref{SIA 262:2025, 3.1.2.2.2}
\begin{equation*}
f_{cm} = f_{ck} + 8\,\mathrm{N}/\mathrm{mm}^{2}
\end{equation*}
\noindent\textbf{Mittelwert des Elastizitätsmoduls} \hfill \normref{SIA 262:2025, 3.1.2.3.3}
\begin{equation*}
E_{cm} = k_e \cdot \sqrt[3]{f_{cm}} \quad \left(f_{cm}\ \text{in}\ \mathrm{N}/\mathrm{mm}^{2}\right)
\end{equation*}
\noindent\textbf{Bemessungswert des Elastizitätsmoduls} \hfill \normref{SIA 262:2025, 4.2.1.15}
\begin{equation*}
E_{cd} = \frac{E_{cm}}{\gamma_{cE}}
\end{equation*}
\noindent\textbf{Krümmungsbeiwert der Spannungs-Dehnungs-Beziehung} \hfill \normref{SIA 262:2025, 4.2.1.6}
\begin{equation*}
k_{\sigma} = \frac{E_{cd}}{400 \cdot f_{cd}}
\end{equation*}
\subsection{Betonstahl}
\noindent\textbf{Bemessungswert der Fliessgrenze} \hfill \normref{SIA 262:2025, 2.4.2.5}
\begin{equation*}
f_{yd} = \frac{f_{yk}}{\gamma_s}
\end{equation*}
\noindent\textbf{Bemessungswert der Fliessgrenze auf Druck} \hfill \normref{SIA 262:2025, 2.4.2.5}
\begin{equation*}
f_{yd}^{-} = \frac{f_{yk}^{-}}{\gamma_s}
\end{equation*}
\subsection{Schiefe Biegung}
Der Querschnitt bleibt eben (Bernoulli). Die Dehnung hängt nur vom Abstand v in Zugrichtung n ab, gemessen ab dem Schwerpunkt; die Nulllinie steht senkrecht zu n und darf schräg liegen. Zu jeder Dehnungsebene ergeben sich N, M\_y und M\_z durch Integration über die Fläche -- als Summe über Streifen quer zu n, deren Grenzen auf jede Polygonecke fallen, und Stab für Stab.

Der Bruchzustand: jede Dehnungsebene, bei der ein Werkstoff an seiner Grenze steht und keiner darüber. Der Beton wird an seinem gedrücktesten Punkt bis -ε\_c2d gestaucht, ganz gedrückt am C-Punkt bis -ε\_c1d; die Bewehrung bis ±ε\_ud. Gesucht wird die Ebene mit N = N\_Ed, deren Moment in die Richtung der Einwirkung zeigt. Nachgewiesen wird nicht die Suche, sondern ihre Probe: die Kräfte jedes Teils und ihre Summen.

\noindent\textbf{Dehnungsebene und innere Kräfte}
\begin{equation*}
\varepsilon(y, z) = \varepsilon_0 + \kappa \cdot v \qquad v = n_y (y - y_S) + n_z (z - z_S) \qquad N = \int_A \sigma\,\mathrm{d}A \qquad M_y = -\int_A \sigma\,(z - z_S)\,\mathrm{d}A \qquad M_z = -\int_A \sigma\,(y - y_S)\,\mathrm{d}A
\end{equation*}
\noindent\textbf{Erfüllungsgrad bei fester Normalkraft und fester Richtung}
\begin{equation*}
\alpha_{eff} = \frac{M_{Rd}}{M_{Ed}} \qquad M = \sqrt{M_y^2 + M_z^2}
\end{equation*}
\noindent\textbf{C30/37: Parabel-Rechteck-Beziehung, Bemessungswerte} \hfill \normref{SIA 262:2025, 4.2.1.6}
\begin{equation*}
\sigma_c(\varepsilon_c) = \begin{cases} -f_{cd} \cdot \dfrac{k_\sigma\,\eta - \eta^2}{1 + (k_\sigma - 2)\,\eta} & 0 \le |\varepsilon_c| \le \varepsilon_{c1d},\ \eta = \dfrac{|\varepsilon_c|}{\varepsilon_{c1d}} \\[2ex] -f_{cd} & \varepsilon_{c1d} < |\varepsilon_c| \le \varepsilon_{c2d} \\[1ex] 0 & \varepsilon_c > 0 \quad (\text{Zug, gerissen}) \end{cases}
\end{equation*}
\noindent\textbf{B500B: bilineare Beziehung, Bemessungswerte} \hfill \normref{SIA 262:2025, 4.2.2.4}
\begin{equation*}
\sigma_s(\varepsilon_s) = \begin{cases} \min(E_s\,\varepsilon_s;\ f_{yd}) & \varepsilon_s \ge 0 \\[1ex] \max(E_s\,\varepsilon_s;\ -f_{yd}^{-}) & \varepsilon_s < 0 \end{cases} \qquad |\varepsilon_s| \le \varepsilon_{ud}
\end{equation*}
\noindent\textbf{Probe: Normalkraft}
\begin{equation*}
N_{Rd} = \sum F_i
\end{equation*}
\noindent\textbf{Probe: Moment um y}
\begin{equation*}
M_{y,Rd} = -\sum F_i \cdot (z_i - z_S)
\end{equation*}
\noindent\textbf{Probe: Moment um z}
\begin{equation*}
M_{z,Rd} = -\sum F_i \cdot (y_i - y_S)
\end{equation*}
\noindent\textbf{Widerstand in der Richtung der Einwirkung}
\begin{equation*}
M_{Rd} = \sqrt{M_{y,Rd}^{2} + M_{z,Rd}^{2}}
\end{equation*}
\noindent\textbf{Erfüllungsgrad}
\begin{equation*}
\alpha_{eff} = \frac{M_{Rd}}{M_{Ed}}
\end{equation*}
\noindent\textbf{C30/37: Spannungsblock 0.85·x, Bemessungswerte}
\begin{equation*}
\sigma_c(\varepsilon_c) = \begin{cases} 0 & \varepsilon_c > -(1 - 0.85) \cdot \varepsilon_{c2d} \\[1ex] -f_{cd} & -\varepsilon_{c2d} \le \varepsilon_c \le -(1 - 0.85) \cdot \varepsilon_{c2d} \end{cases}
\end{equation*}
\noindent\textbf{C30/37: Parabel-Rechteck-Beziehung, charakteristische Werte} \hfill \normref{SIA 262:2025, 4.2.1.6}
\begin{equation*}
\sigma_c(\varepsilon_c) = \begin{cases} -f_{ck} \cdot \dfrac{k_\sigma\,\eta - \eta^2}{1 + (k_\sigma - 2)\,\eta} & 0 \le |\varepsilon_c| \le \varepsilon_{c1d},\ \eta = \dfrac{|\varepsilon_c|}{\varepsilon_{c1d}} \\[2ex] -f_{ck} & \varepsilon_{c1d} < |\varepsilon_c| \le \varepsilon_{c2d} \\[1ex] 0 & \varepsilon_c > 0 \quad (\text{Zug, gerissen}) \end{cases}
\end{equation*}
\noindent\textbf{B500B: bilineare Beziehung, charakteristische Werte} \hfill \normref{SIA 262:2025, 4.2.2.4}
\begin{equation*}
\sigma_s(\varepsilon_s) = \begin{cases} \min(E_s\,\varepsilon_s;\ f_{yk}) & \varepsilon_s \ge 0 \\[1ex] \max(E_s\,\varepsilon_s;\ -f_{yk}^{-}) & \varepsilon_s < 0 \end{cases} \qquad |\varepsilon_s| \le \varepsilon_{ud}
\end{equation*}
\subsection{Querkraft und Torsion}
Jede Schubwand trägt eine Kraft entlang ihrer Achse; ihre Länge ist der Hebelarm des Fachwerks. Verteilt wird elastisch: der Querschnitt verschiebt sich, und jede Wand wehrt sich wie eine Feder mit der Steifigkeit b\_w·l. Das Moment dieser Querkraftanteile um den Schwerpunkt ist der Versatz ΔT; den Rest der Torsion tragen die geschlossenen Zellen als Umlauf-Schubfluss nach Bredt, mehrzellig mit gleicher Verdrillung aller Zellen. Ohne Zelle trägt die Torsion über dieselben Federn mit Drehung. Die Wandkräfte ergeben so immer V\_y, V\_z und um den Schwerpunkt T.

\noindent\textbf{Querkraft über die Federn der Wände}
\begin{equation*}
k_i = b_{w,i} \cdot l_i \qquad \sum_i k_i\, e_i e_i^T\, u = V \qquad V_i = k_i\, e_i \cdot u \qquad \Delta T = \sum_i V_i\, r_i
\end{equation*}
\noindent\textbf{Torsion in den Zellen (Bredt, mehrzellig)}
\begin{equation*}
T - \Delta T = \sum_k 2 A_k\, q_k \qquad \oint_k \frac{q}{b_w}\, \mathrm{d}s = 2 A_k\, \theta
\end{equation*}
\noindent\textbf{Widerstand je Länge einer Wand (Fachwerk)} \hfill \normref{SIA 262:2025, 4.3.3.4}
\begin{equation*}
v_{Rd,s} = \frac{A_{sw}}{s}\, f_{sd} \cot\alpha \qquad v_{Rd,c} = b_w\, k_c\, f_{cd} \sin\alpha \cos\alpha \qquad v_{Rd} = \max_\alpha \min(v_{Rd,s};\ v_{Rd,c})
\end{equation*}
\noindent\textbf{Versatz: Moment der Querkraftanteile um den Schwerpunkt}
\begin{equation*}
\Delta T = \sum_i V_i \cdot r_i
\end{equation*}
\noindent\textbf{Was die Zellen tragen}
\begin{equation*}
T_{Zellen} = T_{Ed} - \Delta T
\end{equation*}
\subsection{Duktilität je Richtung}
Ein Querschnitt kündigt sein Versagen an, wenn die Bewehrung fliesst, bevor der Beton bricht -- dann ist die Druckzone klein. Gemessen wird das an x/d bei reiner Biegung: x ist der Abstand der Nulllinie vom gedrücktesten Betonpunkt, d der des entferntesten gezogenen Stabs, beide senkrecht zur Nulllinie.

\noindent\textbf{Duktilität} \hfill \normref{SIA 262:2025, 4.1.4.2.5}
\begin{equation*}
\frac{x}{d} \le \left(\frac{x}{d}\right)_{max}
\end{equation*}
\noindent\textbf{Bezogene Druckzonenhöhe bei reiner Biegung}
\begin{equation*}
x/d = \frac{x}{d}
\end{equation*}
\subsection{Sprödes Versagen je Richtung}
Ein zu schwach bewehrter Querschnitt reisst und versagt im selben Augenblick. Darum muss der bewehrte Querschnitt mehr tragen als der unbewehrte im Augenblick des Risses: M\_Rd bei N = 0 mindestens das Rissmoment. Das Rissmoment gehört dem ungerissenen Bruttoquerschnitt; bei unsymmetrischer Form mit dem Deviationsmoment.

\noindent\textbf{Spannung im ungerissenen Querschnitt}
\begin{equation*}
\sigma(y, z) = -\frac{M_y\,(I_z z' - I_{yz} y') + M_z\,(I_y y' - I_{yz} z')}{I_y I_z - I_{yz}^2}\qquad y' = y - y_S,\ z' = z - z_S
\end{equation*}
\noindent\textbf{Sprödes Versagen} \hfill \normref{SIA 262:2025, 4.4.1.3}
\begin{equation*}
M_{Rd}(N = 0) \ge M_{Riss} = f_{ct,eff} \cdot W
\end{equation*}
\noindent\textbf{Beiwert für die Höhe in Biegerichtung} \hfill \normref{SIA 262:2025, 4.4.1.3}
\begin{equation*}
k_t = \frac{1}{1 + 0.5 \cdot h/3} \quad \left(h\ \text{in}\ \mathrm{m}\right)
\end{equation*}
\noindent\textbf{Wirksame Zugfestigkeit}
\begin{equation*}
f_{ct,eff} = k_t \cdot f_{ctm}
\end{equation*}
\noindent\textbf{Widerstandsmoment zum gezogensten Betonpunkt}
\begin{equation*}
W = \frac{I_y \cdot I_z - I_{yz}^{2}}{I_{yz} \cdot y' - I_z \cdot z'}
\end{equation*}
\noindent\textbf{Rissmoment des ungerissenen Querschnitts}
\begin{equation*}
M_{Riss} = f_{ct,eff} \cdot W
\end{equation*}
\noindent\textbf{Widerstandsmoment zum gezogensten Betonpunkt}
\begin{equation*}
W = \frac{I_y \cdot I_z - I_{yz}^{2}}{I_z \cdot z' - I_{yz} \cdot y'}
\end{equation*}
\noindent\textbf{Widerstandsmoment zum gezogensten Betonpunkt}
\begin{equation*}
W = \frac{I_y \cdot I_z - I_{yz}^{2}}{I_{yz} \cdot z' - I_y \cdot y'}
\end{equation*}
\noindent\textbf{Widerstandsmoment zum gezogensten Betonpunkt}
\begin{equation*}
W = \frac{I_y \cdot I_z - I_{yz}^{2}}{I_y \cdot y' - I_{yz} \cdot z'}
\end{equation*}
\section{Nachweise}
\subsection{Unterzug}
\noindent\textbf{Angaben zum Querschnitt}
\begin{equation*}
\text{C30/37} \qquad A = 180000\,\mathrm{mm}^2 \qquad y_S = 150.0\,\mathrm{mm} \qquad z_S = 300.0\,\mathrm{mm}
\end{equation*}
\noindent\textbf{Bewehrung und Schubwände}
\begin{longtable}{lll}
\hline
Element & Bewehrung & Stahl \\
\hline
\endhead
Linie 1 & $3 \varnothing 20$ & B500B \\
Linie 2 & $2 \varnothing 12$ & B500B \\
Wand 1 & $b_w = 100,\ 1 \times \varnothing 10\ @\ 150$ & B500B \\
Wand 2 & $b_w = 100,\ 1 \times \varnothing 10\ @\ 150$ & B500B \\
Wand 3 & $b_w = 100,\ 1 \times \varnothing 10\ @\ 150$ & B500B \\
Wand 4 & $b_w = 100,\ 1 \times \varnothing 10\ @\ 150$ & B500B \\
\hline
\end{longtable}
\begin{xltabular}{\linewidth}{>{\raggedright\arraybackslash}X>{\raggedright\arraybackslash}Xrrr}
\hline
Nachweis & Bezeichnung & Widerstand & Einwirkung & $\alpha_{eff}$ \\
\hline
\endhead
Biegung und Normalkraft & Feld & $M_{Rd} = 211.8\,\mathrm{kNm}$ & $M_{Ed} = 150.0\,\mathrm{kNm}$ & $1.41$ \\
Querkraft und Torsion & Feld & $v_{Rd} = 394.3\,\mathrm{kN}/\mathrm{m}$ & $q_{Ed} = 225.0\,\mathrm{kN}/\mathrm{m}$ & $1.75$ \\
\hline
\end{xltabular}
\subsection{Kastenträger}
\noindent\textbf{Angaben zum Querschnitt}
\begin{equation*}
\text{C30/37} \qquad A = 330000\,\mathrm{mm}^2 \qquad y_S = 400.0\,\mathrm{mm} \qquad z_S = 300.0\,\mathrm{mm}
\end{equation*}
\noindent\textbf{Bewehrung und Schubwände}
\begin{longtable}{lll}
\hline
Element & Bewehrung & Stahl \\
\hline
\endhead
Stab 1 & $\varnothing 20$ & B500B \\
Linie 1 & $6 \varnothing 16$ & B500B \\
Linie 2 & $4 \varnothing 12$ & B500B \\
Linie 3 & $400\,\mathrm{mm}^2$ & B500B \\
Wand 1 & $b_w = 150,\ 1 \times \varnothing 10\ @\ 150$ & B500B \\
Wand 2 & $b_w = 150,\ 1 \times \varnothing 10\ @\ 150$ & B500B \\
Wand 3 & $b_w = 150,\ 1 \times \varnothing 10\ @\ 150$ & B500B \\
Wand 4 & $b_w = 150,\ 1 \times \varnothing 10\ @\ 150$ & B500B \\
\hline
\end{longtable}
\begin{xltabular}{\linewidth}{>{\raggedright\arraybackslash}X>{\raggedright\arraybackslash}Xrrr}
\hline
Nachweis & Bezeichnung & Widerstand & Einwirkung & $\alpha_{eff}$ \\
\hline
\endhead
Querkraft und Torsion & Feld & $v_{Rd} = 394.3\,\mathrm{kN}/\mathrm{m}$ & $q_{Ed} = 380.3\,\mathrm{kN}/\mathrm{m}$ & $1.04$ \\
Querkraft und Torsion & Schief & $v_{Rd} = 394.3\,\mathrm{kN}/\mathrm{m}$ & $q_{Ed} = 76.9\,\mathrm{kN}/\mathrm{m}$ & $5.13$ \\
Biegung und Normalkraft & Feld & $M_{Rd} = 188.6\,\mathrm{kNm}$ & $M_{Ed} = 304.5\,\mathrm{kNm}$ & $0.62$ \\
Biegung und Normalkraft & Schief & $M_{Rd} = 666.6\,\mathrm{kNm}$ & $M_{Ed} = 291.5\,\mathrm{kNm}$ & $2.29$ \\
Duktilität & Zug links & $\left(x/d\right)_{max} = 0.35$ & $x/d = 0.135$ & $2.59$ \\
Sprödes Versagen & Zug oben & $M_{Rd} = 194.7\,\mathrm{kNm}$ & $M_{Riss} = 116.7\,\mathrm{kNm}$ & $1.67$ \\
\hline
\end{xltabular}
\hinweis{Warnung}{Biegung und Normalkraft – Feld: nicht erfüllt. Bei N\_Ed = 664.0 kN, in der Richtung des Moments.}
\subsection{Stütze}
\noindent\textbf{Angaben zum Querschnitt}
\begin{equation*}
\text{C30/37} \qquad A = 160000\,\mathrm{mm}^2 \qquad y_S = 200.0\,\mathrm{mm} \qquad z_S = 200.0\,\mathrm{mm}
\end{equation*}
\noindent\textbf{Bewehrung und Schubwände}
\begin{longtable}{lll}
\hline
Element & Bewehrung & Stahl \\
\hline
\endhead
Linie 1 & $3 \varnothing 20$ & B500B \\
Linie 2 & $3 \varnothing 20$ & B500B \\
\hline
\end{longtable}
\begin{xltabular}{\linewidth}{>{\raggedright\arraybackslash}X>{\raggedright\arraybackslash}Xrrr}
\hline
Nachweis & Bezeichnung & Widerstand & Einwirkung & $\alpha_{eff}$ \\
\hline
\endhead
Biegung und Normalkraft & Druck & $M_{Rd} = 348.9\,\mathrm{kNm}$ & $M_{Ed} = 120.0\,\mathrm{kNm}$ & $2.91$ \\
Biegung und Normalkraft & Zug & $N_{Rd}^{+} = 942.5\,\mathrm{kN}$ & $N_{Ed} = 400.0\,\mathrm{kN}$ & $2.36$ \\
\hline
\end{xltabular}
Mindestens ein Nachweis ist nicht erfüllt.

\section{Werte}
\begin{xltabular}{\linewidth}{>{\raggedright\arraybackslash}Xlrll}
\hline
Bezeichnung & Symbol & Wert & Einheit & Herkunft \\
\hline
\endhead
Bemessungswert des Elastizitätsmoduls & $E_{cd}$ & $33620$ & N/mm² & berechnet \\
Mittelwert des Elastizitätsmoduls & $E_{cm}$ & $33620$ & N/mm² & berechnet \\
Dehnung am Ende des ansteigenden Astes & $\varepsilon_{c1d}$ & $2$ & ‰ & Vorgabe \\
Bruchdehnung des Betons & $\varepsilon_{c2d}$ & $3.5$ & ‰ & Vorgabe \\
Beiwert zur Berücksichtigung der Festigkeitsminderung & $\eta_{fc}$ & $1$ &  & berechnet \\
Bemessungswert der Betondruckfestigkeit & $f_{cd}$ & $20$ & N/mm² & berechnet \\
Charakteristische Zylinderdruckfestigkeit & $f_{ck}$ & $30$ & N/mm² & Vorgabe \\
Mittelwert der Zylinderdruckfestigkeit & $f_{cm}$ & $38$ & N/mm² & berechnet \\
Mittelwert der Zugfestigkeit & $f_{ctm}$ & $2.9$ & N/mm² & Vorgabe \\
Teilsicherheitsbeiwert für Beton & $\gamma_c$ & $1.5$ &  & Vorgabe \\
Teilsicherheitsbeiwert für den Elastizitätsmodul & $\gamma_{cE}$ & $1$ &  & Vorgabe \\
Beiwert für die Gesteinskörnung & $k_e$ & $10000$ &  & Vorgabe \\
Krümmungsbeiwert der Spannungs-Dehnungs-Beziehung & $k_{\sigma}$ & $4.202$ &  & berechnet \\
Elastizitätsmodul des Betonstahls & $E_s$ & $200000$ & N/mm² & Vorgabe \\
Bemessungswert der Dehnung bei Höchstlast & $\varepsilon_{ud}$ & $4.5$ & \% & Vorgabe \\
Bemessungswert der Fliessgrenze & $f_{yd}$ & $435$ & N/mm² & berechnet \\
Bemessungswert der Fliessgrenze auf Druck & $f_{yd}^{-}$ & $435$ & N/mm² & berechnet \\
Charakteristische Fliessgrenze & $f_{yk}$ & $500$ & N/mm² & Vorgabe \\
Charakteristische Fliessgrenze auf Druck & $f_{yk}^{-}$ & $500$ & N/mm² & Vorgabe \\
Teilsicherheitsbeiwert für Betonstahl & $\gamma_s$ & $1.15$ &  & Vorgabe \\
Fläche des Bruttoquerschnitts & $A$ & $180000$ & mm² & berechnet \\
Stahlfläche der Bewehrung & $A_{s,tot}$ & $1169$ & mm² & berechnet \\
Trägheitsmoment um y & $I_y$ & $5400$ & ·10⁶ mm⁴ & berechnet \\
Deviationsmoment & $I_{yz}$ & $0$ & ·10⁶ mm⁴ & berechnet \\
Trägheitsmoment um z & $I_z$ & $1350$ & ·10⁶ mm⁴ & berechnet \\
Grösste Neigung der Druckdiagonalen & $\alpha_{max}$ & $45$ & ° & Vorgabe \\
Kleinste Neigung der Druckdiagonalen & $\alpha_{min}$ & $30$ & ° & Vorgabe \\
Abminderung der Betondruckfestigkeit in der Druckdiagonalen & $k_c$ & $0.55$ &  & Vorgabe \\
Erfüllungsgrad Duktilität – Zug links & $\alpha_{eff,x/d,\text{Zug links}}$ & $1.60$ &  & berechnet \\
Erfüllungsgrad Duktilität – Zug oben & $\alpha_{eff,x/d,\text{Zug oben}}$ & $4.56$ &  & berechnet \\
Erfüllungsgrad Duktilität – Zug rechts & $\alpha_{eff,x/d,\text{Zug rechts}}$ & $1.60$ &  & berechnet \\
Erfüllungsgrad Duktilität – Zug unten & $\alpha_{eff,x/d,\text{Zug unten}}$ & $2.73$ &  & berechnet \\
Erfüllungsgrad Biegung und Normalkraft – Feld & $\alpha_{eff,\text{Feld}}$ & $1.41$ &  & berechnet \\
Grösste aufnehmbare Druckkraft & $N_{Rd}^{-}$ & $-4044.1$ & kN & berechnet \\
Grösste aufnehmbare Zugkraft & $N_{Rd}^{+}$ & $508.1$ & kN & berechnet \\
Erfüllungsgrad sprödes Versagen – Zug links & $\alpha_{eff,Riss,\text{Zug links}}$ & $2.45$ &  & berechnet \\
Erfüllungsgrad sprödes Versagen – Zug oben & $\alpha_{eff,Riss,\text{Zug oben}}$ & $1.18$ &  & berechnet \\
Erfüllungsgrad sprödes Versagen – Zug rechts & $\alpha_{eff,Riss,\text{Zug rechts}}$ & $2.45$ &  & berechnet \\
Erfüllungsgrad sprödes Versagen – Zug unten & $\alpha_{eff,Riss,\text{Zug unten}}$ & $4.46$ &  & berechnet \\
Erfüllungsgrad Querkraft und Torsion – Feld & $\alpha_{eff,V,\text{Feld}}$ & $1.75$ &  & berechnet \\
Schwerpunkt, y & $y_S$ & $150$ & mm & berechnet \\
Schwerpunkt, z & $z_S$ & $300$ & mm & berechnet \\
Fläche des Bruttoquerschnitts & $A$ & $330000$ & mm² & berechnet \\
Stahlfläche der Bewehrung & $A_{s,tot}$ & $2373$ & mm² & berechnet \\
Trägheitsmoment um y & $I_y$ & $13275$ & ·10⁶ mm⁴ & berechnet \\
Deviationsmoment & $I_{yz}$ & $0$ & ·10⁶ mm⁴ & berechnet \\
Trägheitsmoment um z & $I_z$ & $22475$ & ·10⁶ mm⁴ & berechnet \\
Grösste Neigung der Druckdiagonalen & $\alpha_{max}$ & $45$ & ° & Vorgabe \\
Kleinste Neigung der Druckdiagonalen & $\alpha_{min}$ & $30$ & ° & Vorgabe \\
Abminderung der Betondruckfestigkeit in der Druckdiagonalen & $k_c$ & $0.55$ &  & Vorgabe \\
Erfüllungsgrad Duktilität – Zug links & $\alpha_{eff,x/d,\text{Zug links}}$ & $2.59$ &  & berechnet \\
Erfüllungsgrad Duktilität – Zug oben & $\alpha_{eff,x/d,\text{Zug oben}}$ & $4.09$ &  & berechnet \\
Erfüllungsgrad Duktilität – Zug rechts & $\alpha_{eff,x/d,\text{Zug rechts}}$ & $2.67$ &  & berechnet \\
Erfüllungsgrad Duktilität – Zug unten & $\alpha_{eff,x/d,\text{Zug unten}}$ & $3.02$ &  & berechnet \\
Erfüllungsgrad Biegung und Normalkraft – Feld & $\alpha_{eff,\text{Feld}}$ & $0.62$ &  & berechnet \\
Grösste aufnehmbare Druckkraft & $N_{Rd}^{-}$ & $-7501.7$ & kN & berechnet \\
Grösste aufnehmbare Zugkraft & $N_{Rd}^{+}$ & $1031.7$ & kN & berechnet \\
Erfüllungsgrad Biegung und Normalkraft – Schief & $\alpha_{eff,\text{Schief}}$ & $2.29$ &  & berechnet \\
Erfüllungsgrad sprödes Versagen – Zug links & $\alpha_{eff,Riss,\text{Zug links}}$ & $2.70$ &  & berechnet \\
Erfüllungsgrad sprödes Versagen – Zug oben & $\alpha_{eff,Riss,\text{Zug oben}}$ & $1.67$ &  & berechnet \\
Erfüllungsgrad sprödes Versagen – Zug rechts & $\alpha_{eff,Riss,\text{Zug rechts}}$ & $2.52$ &  & berechnet \\
Erfüllungsgrad sprödes Versagen – Zug unten & $\alpha_{eff,Riss,\text{Zug unten}}$ & $3.08$ &  & berechnet \\
Moment der Längszugkraft um y – Feld & $\Delta M_y$ & $0$ & kNm & berechnet \\
Moment der Längszugkraft um z – Feld & $\Delta M_z$ & $-52$ & kNm & berechnet \\
Längszugkraft aus Querkraft und Torsion – Feld & $\Delta N$ & $664$ & kN & berechnet \\
Erfüllungsgrad Querkraft und Torsion – Feld & $\alpha_{eff,V,\text{Feld}}$ & $1.04$ &  & berechnet \\
Moment der Längszugkraft um y – Schief & $\Delta M_y$ & $0$ & kNm & berechnet \\
Moment der Längszugkraft um z – Schief & $\Delta M_z$ & $0$ & kNm & berechnet \\
Längszugkraft aus Querkraft und Torsion – Schief & $\Delta N$ & $173.2$ & kN & berechnet \\
Erfüllungsgrad Querkraft und Torsion – Schief & $\alpha_{eff,V,\text{Schief}}$ & $5.13$ &  & berechnet \\
Schwerpunkt, y & $y_S$ & $400$ & mm & berechnet \\
Schwerpunkt, z & $z_S$ & $300$ & mm & berechnet \\
Fläche des Bruttoquerschnitts & $A$ & $160000$ & mm² & berechnet \\
Stahlfläche der Bewehrung & $A_{s,tot}$ & $1885$ & mm² & berechnet \\
Trägheitsmoment um y & $I_y$ & $2133.3$ & ·10⁶ mm⁴ & berechnet \\
Deviationsmoment & $I_{yz}$ & $0$ & ·10⁶ mm⁴ & berechnet \\
Trägheitsmoment um z & $I_z$ & $2133.3$ & ·10⁶ mm⁴ & berechnet \\
Erfüllungsgrad Duktilität – Zug oben & $\alpha_{eff,x/d,\text{Zug oben}}$ & $2.52$ &  & berechnet \\
Erfüllungsgrad Duktilität – Zug unten & $\alpha_{eff,x/d,\text{Zug unten}}$ & $2.52$ &  & berechnet \\
Erfüllungsgrad Biegung und Normalkraft – Druck & $\alpha_{eff,\text{Druck}}$ & $2.91$ &  & berechnet \\
Grösste aufnehmbare Druckkraft & $N_{Rd}^{-}$ & $-5497.4$ & kN & berechnet \\
Grösste aufnehmbare Zugkraft & $N_{Rd}^{+}$ & $942.5$ & kN & berechnet \\
Erfüllungsgrad Biegung und Normalkraft – Zug & $\alpha_{eff,\text{Zug}}$ & $2.36$ &  & berechnet \\
Erfüllungsgrad sprödes Versagen – Zug oben & $\alpha_{eff,Riss,\text{Zug oben}}$ & $5.37$ &  & berechnet \\
Erfüllungsgrad sprödes Versagen – Zug unten & $\alpha_{eff,Riss,\text{Zug unten}}$ & $5.37$ &  & berechnet \\
Schwerpunkt, y & $y_S$ & $200$ & mm & berechnet \\
Schwerpunkt, z & $z_S$ & $200$ & mm & berechnet \\
\hline
\end{xltabular}
\end{document}
